% ==========================================================================
%  CODE_EXPLANATION.tex  ---  Comprehensive walkthrough of the GA Prompt Optimiser
% ==========================================================================
\documentclass[11pt,a4paper]{report}

% --- Encoding & fonts ---
\usepackage[utf8]{inputenc}
\usepackage[T1]{fontenc}
\usepackage{lmodern}

% --- Layout ---
\usepackage[margin=2.4cm,top=2.8cm,bottom=2.8cm]{geometry}
\usepackage{fancyhdr}
\usepackage{titlesec}

% --- Colour ---
\usepackage{xcolor}
\definecolor{codebg}{HTML}{F7F7F7}
\definecolor{codeframe}{HTML}{CCCCCC}
\definecolor{kwcolor}{HTML}{0000CC}
\definecolor{strcolor}{HTML}{007700}
\definecolor{cmtcolor}{HTML}{888888}
\definecolor{conceptbg}{HTML}{EBF5FB}
\definecolor{conceptframe}{HTML}{2980B9}
\definecolor{warnbg}{HTML}{FFF3E0}
\definecolor{warnframe}{HTML}{E65100}
\definecolor{specbg}{HTML}{F0FFF0}
\definecolor{specframe}{HTML}{27AE60}
\definecolor{chainbg}{HTML}{FAFAFA}
\definecolor{chainframe}{HTML}{7F8C8D}
\definecolor{exbg}{HTML}{FFFDE7}
\definecolor{exframe}{HTML}{F39C12}

% --- Code listings ---
\usepackage{listings}
\lstdefinestyle{pycode}{
  language=Python,
  basicstyle=\ttfamily\small,
  keywordstyle=\color{kwcolor}\bfseries,
  stringstyle=\color{strcolor},
  commentstyle=\color{cmtcolor}\itshape,
  backgroundcolor=\color{codebg},
  frame=single,
  rulecolor=\color{codeframe},
  breaklines=true,
  showstringspaces=false,
  numbers=left,
  numberstyle=\tiny\color{gray},
  numbersep=8pt,
  tabsize=4,
  xleftmargin=1.5em,
  framexleftmargin=1.5em,
  aboveskip=0.8em,
  belowskip=0.8em,
  literate={---}{{---}}1,
}
\lstdefinestyle{chain}{
  language={},
  basicstyle=\ttfamily\footnotesize,
  backgroundcolor=\color{chainbg},
  frame=single,
  rulecolor=\color{chainframe},
  breaklines=true,
  showstringspaces=false,
  numbers=none,
  aboveskip=0.6em,
  belowskip=0.6em,
  xleftmargin=0.5em,
  framexleftmargin=0.5em,
}
\lstdefinestyle{example}{
  language=Python,
  basicstyle=\ttfamily\footnotesize,
  keywordstyle=\color{kwcolor}\bfseries,
  stringstyle=\color{strcolor},
  commentstyle=\color{cmtcolor}\itshape,
  backgroundcolor=\color{exbg},
  frame=single,
  rulecolor=\color{exframe},
  breaklines=true,
  showstringspaces=false,
  numbers=none,
  aboveskip=0.6em,
  belowskip=0.6em,
  xleftmargin=0.5em,
  framexleftmargin=0.5em,
}
\lstset{style=pycode}

% --- Boxes ---
\usepackage{tcolorbox}
\tcbuselibrary{skins,breakable}

\newtcolorbox{conceptbox}[1]{
  colback=conceptbg, colframe=conceptframe,
  fonttitle=\bfseries, title={#1},
  boxrule=0.8pt, arc=3pt, left=6pt, right=6pt, top=4pt, bottom=4pt,
  breakable,
}
\newtcolorbox{warnbox}[1]{
  colback=warnbg, colframe=warnframe,
  fonttitle=\bfseries, title={#1},
  boxrule=0.8pt, arc=3pt, left=6pt, right=6pt, top=4pt, bottom=4pt,
  breakable,
}
\newtcolorbox{specbox}[1]{
  colback=specbg, colframe=specframe,
  fonttitle=\bfseries\ttfamily, title={#1},
  boxrule=0.8pt, arc=3pt, left=6pt, right=6pt, top=4pt, bottom=4pt,
  breakable,
}
\newtcolorbox{exbox}[1]{
  colback=exbg, colframe=exframe,
  fonttitle=\bfseries, title={Worked Example --- #1},
  boxrule=0.8pt, arc=3pt, left=6pt, right=6pt, top=4pt, bottom=4pt,
  breakable,
}

% --- Tables ---
\usepackage{booktabs}
\usepackage{longtable}
\usepackage{array}
\usepackage{tabularx}

% --- Lists ---
\usepackage{enumitem}

% --- Graphics ---
\usepackage{graphicx}

% --- Hyperref ---
\usepackage{hyperref}
\hypersetup{
  colorlinks=true,
  linkcolor=blue!60!black,
  urlcolor=blue!70!black,
  pdftitle={Code Explanation -- GA Prompt Optimiser},
  pdfauthor={Group 1},
}

% --- Headers ---
\pagestyle{fancy}
\fancyhf{}
\fancyhead[L]{\small\leftmark}
\fancyhead[R]{\small Code Explanation}
\fancyfoot[C]{\thepage}
\renewcommand{\headrulewidth}{0.4pt}

% --- Chapter style ---
\titleformat{\chapter}[display]
  {\normalfont\huge\bfseries}{\chaptertitlename\ \thechapter}{18pt}{\Huge}
\titlespacing*{\chapter}{0pt}{-20pt}{30pt}

% --- Convenience macros ---
\newcommand{\pyfile}[1]{\texttt{#1}}
\newcommand{\fn}[1]{\texttt{#1}}
\newcommand{\param}[1]{\textit{#1}}

% ==========================================================================
\begin{document}

% ---------------------------------------------------------------------------
%  TITLE PAGE
% ---------------------------------------------------------------------------
\begin{titlepage}
\centering
\vspace*{2.5cm}
{\Huge\bfseries Code Explanation\par}
\vspace{0.8cm}
{\LARGE GA Prompt Optimiser for LLM Code Generation\par}
\vspace{1.2cm}
{\Large Evolving System Prompts on Codeforces Problems\par}
\vspace{2cm}
{\large Group 1\par}
\vspace{0.4cm}
{\large Computational Intelligence --- Spring 2026\par}
\vspace{1cm}
{\large April 2026\par}
\vfill
{\small
  Deep-dive walkthrough of every file, every function, and every call chain.\\[4pt]
  Includes exact function signatures, cross-file data-flow diagrams, and\\
  step-by-step worked examples tracing real data through the pipeline.\\[4pt]
  Backend: Google Cloud Vertex AI (Gemini 2.5 Flash \& Pro).\\
  Designed for readers with no prior knowledge of genetic algorithms or LLM APIs.
}
\end{titlepage}

\tableofcontents
\newpage

% ==========================================================================
%  CHAPTER 1 --- INTRODUCTION
% ==========================================================================
\chapter{Introduction}

\section{What This Project Does}

This project uses a \textbf{Genetic Algorithm (GA)} to automatically discover the best
``system prompt'' for a Large Language Model (LLM) that generates \textbf{C++ code}
for competitive programming problems.

A \textbf{system prompt} is a paragraph of plain-English instructions prepended to every
LLM call, shaping how it thinks before it sees the coding problem.
Examples: \emph{``Think like a competitive programmer''} or \emph{``Always check edge
cases before writing code.''}

The GA treats each system prompt as a \textbf{chromosome} --- a candidate solution that
can be scored (fitness evaluated), combined with another (crossover), and lightly altered
(mutation), exactly as in biological evolution.  Over 12 generations the GA evolves a
population of 5 prompts, keeping the best-performing ones and combining their strategies,
until it converges on a prompt that makes the LLM solve the most Codeforces problems correctly.

\begin{conceptbox}{The Core Insight}
Manual prompt engineering is trial-and-error.  A human tries one prompt, tests it, tweaks
it, and tries again.  This project automates that loop: the GA \emph{searches} the entire
space of possible English instruction paragraphs by treating test-case pass rates as the
fitness signal.  The LLM's own creativity (via crossover and mutation calls) is used to
\emph{generate} new prompts, while its code-generation ability is what is being
\emph{optimised}.
\end{conceptbox}

\begin{conceptbox}{Key design choice: C++ and Vertex AI}
The target model is asked to produce a \textbf{complete C++ program} (\fn{int main()},
reads from \fn{stdin}, writes to \fn{stdout}).  Evaluation compiles the code with
\fn{g++ -O2 -std=c++23} and runs the resulting binary.  All LLM calls go through the
\textbf{Google Cloud Vertex AI} REST API via the \texttt{google-genai} Python SDK,
authenticated with a service account JSON file --- no subscription or interactive login
is required.
\end{conceptbox}

\section{The Codebase at a Glance}

The project consists of 9 Python files:

\begin{table}[h]
\centering
\begin{tabular}{@{}llrp{5.5cm}@{}}
\toprule
\textbf{File} & \textbf{Lines} & \textbf{Public items} & \textbf{Role} \\
\midrule
\pyfile{config.py}        & 83  & 25 constants  & Single source of truth for all settings \\
\pyfile{seeds.py}         & 137 & 1 list        & The 15 generation-zero C++ chromosomes \\
\pyfile{dataset.py}       & 253 & 3 functions   & Load \& stratified-split 305 Codeforces problems \\
\pyfile{evaluator.py}     & 295 & 4 functions   & Extract C++ code; compile; run binary \\
\pyfile{cost\_tracker.py} & 244 & 2 classes     & Stream token counts; compute cost; write JSON \\
\pyfile{llm.py}           & 470 & 7 functions   & All Vertex AI Gemini calls via google-genai SDK \\
\pyfile{ga.py}            & 717 & 10 items      & The full evolutionary loop (parallel eval, resume, selection schemes) \\
\pyfile{main.py}          & 545 & 6 functions   & Orchestrate, checkpoint, resume, save results \\
\pyfile{genealogy.py}     & 154 & 1 function    & Ancestry report for the best evolved prompt \\
\bottomrule
\end{tabular}
\caption{Project file summary.}
\end{table}

\section{Module Dependency Graph}

\begin{lstlisting}[style=chain, caption={Module import graph (arrows = imports).}]
  main.py
    |--- imports ---> config, dataset, evaluator, ga, llm, genealogy
    |--- imports ---> cost_tracker  (singleton: tracker)
    |--- imports ---> seeds         (SEED_PROMPTS list)

  ga.py
    |--- imports ---> config, evaluator, llm

  llm.py
    |--- imports ---> config
    |--- imports ---> cost_tracker  (singleton: tracker, PRICING dict)
    |--- uses     ---> google.genai SDK  (Vertex AI realtime API)
    |--- uses     ---> google.oauth2.service_account  (auth)

  evaluator.py
    |--- imports ---> os, re, signal, subprocess, tempfile, time  (stdlib only)
    |--- imports ---> config

  dataset.py
    |--- imports ---> json, random, sys, pathlib, collections  (stdlib only)
    |--- imports ---> config

  cost_tracker.py
    |--- imports ---> json, threading, time, dataclasses, datetime, pathlib  (stdlib)

  genealogy.py
    |--- imports ---> json, pathlib  (stdlib only)

  config.py
    |--- imports ---> pathlib  (stdlib only)

  seeds.py
    |--- (no imports)
\end{lstlisting}

\begin{conceptbox}{Key Import Insight}
\pyfile{evaluator.py}, \pyfile{dataset.py}, \pyfile{cost\_tracker.py}, \pyfile{genealogy.py},
\pyfile{config.py}, and \pyfile{seeds.py} have \emph{zero external dependencies}.
\pyfile{llm.py} is the \emph{only} file that touches the network, and it does so by calling
the \texttt{google-genai} Python SDK which connects to Google Cloud Vertex AI via HTTPS,
authenticated with a service account JSON credential file.
\end{conceptbox}

\section{How to Read This Document}

Each chapter covers one file and follows this structure:
\begin{enumerate}[nosep]
  \item \textbf{Purpose} --- plain-English summary of what the file does and why it exists.
  \item \textbf{File at a Glance} --- table of every function/class with one-line descriptions.
  \item \textbf{Full Source Code} --- the actual source with line numbers.
  \item \textbf{Function-by-Function Analysis} --- for every function:
    \begin{itemize}[nosep]
      \item Exact signature (all parameters, types, defaults)
      \item What calls it and what it calls (the call chain)
      \item Step-by-step walkthrough of the body
      \item Return value with a concrete example
      \item Edge cases and error handling
    \end{itemize}
  \item \textbf{Worked Example} --- real data traced through the function.
\end{enumerate}

Colour coding for boxes: \textbf{blue} = concept/theory; \textbf{orange} = warning/gotcha;
\textbf{green} = function specification; \textbf{yellow} = worked example.


% ==========================================================================
%  CHAPTER 2 --- ARCHITECTURE \& DATA FLOW
% ==========================================================================
\chapter{Architecture \& Data Flow}

\section{The Two-Model Architecture}

A critical design choice is that \emph{two separate Gemini models} are used for two
completely different purposes:

\begin{table}[h]
\centering
\begin{tabular}{@{}lllp{5cm}@{}}
\toprule
\textbf{Role} & \textbf{Model} & \textbf{Temperature} & \textbf{Used for} \\
\midrule
\textbf{Target} (the subject)   & \texttt{gemini-2.5-flash} & 0.0 & Generating C++ code (fitness eval) \\
\textbf{Optimizer} (the tool)   & \texttt{gemini-2.5-pro}   & 0.7 & Crossover \& mutation of prompts \\
\bottomrule
\end{tabular}
\caption{The two-model design. Temperature is enforced via the API \texttt{GenerateContentConfig}.}
\end{table}

\begin{conceptbox}{Why Two Models?}
\textbf{Target} (\texttt{gemini-2.5-flash}): We are searching for the prompt that makes
\emph{this specific model} write the best C++ code.  The goal is to fix the model and
optimise the prompt.  Temperature 0.0 gives deterministic, reproducible code output.\\[6pt]
\textbf{Optimizer} (\texttt{gemini-2.5-pro}): Crossover and mutation require
\emph{creativity} --- merging two prompts into a coherent paragraph, inventing a new
instructional strategy.  The Pro model is more capable for creative reasoning tasks.
Temperature 0.7 adds controlled randomness to generate diverse offspring.
\end{conceptbox}

\begin{conceptbox}{Thinking tokens}
Both models are configured with a \texttt{ThinkingConfig}: \texttt{thinking\_budget=512}
for fitness calls and \texttt{thinking\_budget=2048} for operator calls.  Gemini 2.5
Flash/Pro can use ``thinking tokens'' (internal chain-of-thought) before generating the
final output.  These are billed at the output token rate and tracked separately in
\fn{usage\_metadata.thoughts\_token\_count}.
\end{conceptbox}

\section{Master Execution Pipeline}

Below is the complete call tree for a single run.

\begin{lstlisting}[style=chain, caption={Complete call tree for one GA run.}]
main.main()
  |
  |-- [pre-flight: Vertex AI credentials]
  |     check config.SERVICE_ACCOUNT_PATH exists
  |     import google.genai, google.oauth2.service_account  (verify installed)
  |     read service-account.json -> verify project_id, client_email
  |
  |-- [pre-flight: dataset dirs]
  |     config.APPS_TRAIN.exists() and config.APPS_TEST.exists()
  |
  |-- dataset.load_problems()
  |     Reads 305 pre-selected Codeforces problems from _SELECTED dict
  |     Returns: list[dict]  (id, question, inputs, outputs, rating)
  |
  |-- dataset.split(all_problems, 65, 200, seed=42)
  |     Stratified: harder train (>=1800) always in evolution
  |     Regular tiers 800-1700 drawn per _EVOLUTION_PER_TIER counts
  |     Returns: (evolution[65], holdout[200], rest[~40])
  |
  |-- tracker.set_log_path(run_dir/"calls.jsonl")
  |-- tracker.set_live_path(run_dir/"costs_live.json")
  |-- tracker.set_costs_live_path(Costs/"cost_<run>.json")
  |-- llm.set_error_log_path(run_dir/"llm_errors.log")
  |
  +-- ga.run(SEED_PROMPTS[:5], evolution_pool, run_dir, on_generation=...,
             selection_scheme=args.scheme, elitism_count=1 if scheme=="elitism" else None)
        |
        | [GEN 0: evaluate all 5 seeds, sequential]
        |   for seed_idx, ind in enumerate(population):
        |     [restore from seed checkpoint if it exists]
        |     ind.fitness = compute_fitness(gen=0, ind, problems)
        |       ThreadPoolExecutor(max_workers=5):
        |         _eval_one_problem(gen, ind, prob, idx, total)  x65
        |           time.sleep(random.uniform(0, 1))   -- stagger
        |           llm.generate_code(ind.prompt, prob["question"])
        |             client.models.generate_content(
        |                 model="gemini-2.5-flash",
        |                 contents=user_msg,
        |                 config=GenerateContentConfig(
        |                     system_instruction=system,
        |                     max_output_tokens=8192,
        |                     temperature=0.0,
        |                     thinking_config=ThinkingConfig(budget=512)))
        |             tracker.record(phase,model,in,out,cost)
        |             returns text (raw LLM output)
        |           evaluator.extract_code(response) -> C++ string
        |           evaluator.run_tests_verbose(code, inputs, outputs, 2)
        |             _compile(code)  -> g++ -O2 -std=c++23 -> binary
        |             for each test: _run_one(bin, stdin, 2) -> stdout
        |             _outputs_match(actual, expected) -> bool
        |             returns {score, reason, compile, test_cases, ...}
        |           _log_evaluation(entry) -> evaluations_genNN.jsonl
        |           returns (score, failure_example | None)
        |     _assign_failure_examples(ind, collected_failures)
        |     save seed checkpoint: seed_NN_<uid>.json
        |
        | [GENS 1..12]
        |   population.sort(by fitness, descending)
        |   stats = {generation, best, mean, worst, std, cv, best_uid, best_prompt}
        |   history.append(stats)
        |   on_generation(gen-1, population, stats, rng.getstate())
        |     -> save_checkpoint(run_dir, gen-1, pop, stats, rng_state)
        |        writes: results/run_*/generation_NN.json  (includes rng_state)
        |
        |   [convergence check: CV < 0.01 for 3 consecutive gens -> early stop]
        |
        |   new_pop = []   -- ELITISM_COUNT=0, no frozen elites
        |   offspring_idx = 0
        |   while len(new_pop) < 5:   -- breed 5 offspring
        |     p1 = tournament_select(population, rng)  -- best of 3
        |     p2 = tournament_select(population, rng)
        |     allow_mut = (offspring_idx >= CROSSOVER_ONLY_SLOTS=2)
        |     child = breed(p1, p2, rng, allow_mutation=allow_mut)
        |       [80%] llm.crossover(a.prompt, b.prompt)  -> merged prompt
        |       [if allow_mut:
        |         [25%] llm.mutate_inject(prompt, failure_examples)
        |         [15%] llm.mutate_delete(prompt)
        |         [10%] llm.mutate_rephrase(prompt)]
        |       returns Individual(prompt, parent_a, parent_b, operators,
        |                         intermediate_steps)
        |     new_pop.append(child); offspring_idx += 1
        |
        |   for ind in new_pop:   -- evaluate all 5 offspring
        |     ind.fitness = compute_fitness(gen, ind, problems)
        |   population = new_pop
        |
        +-- returns (best_Individual, history[list of dicts])
  |
  |-- main.evaluate_on_holdout(best.prompt, holdout[200])
  |     ThreadPoolExecutor(max_workers=5): 200 calls, same pipeline
  |     Returns: float (holdout fitness)
  |
  |-- genealogy.generate_report(run_dir, best.uid)
  |     reads generation_NN.json files -> builds ancestry chain
  |     writes genealogy_report.txt + genealogy_tree.json
  |
  |-- (run_dir/"summary.json").write_text(...)
  +-- tracker.flush(extra={...})
        writes: Code/Costs/cost_<run>.json
        Returns: Path
\end{lstlisting}

\section{Data Types Flowing Between Modules}

\begin{table}[h]
\centering
\begin{tabular}{@{}p{3cm}p{3.2cm}p{2.5cm}p{5cm}@{}}
\toprule
\textbf{Producer} & \textbf{Consumer} & \textbf{Type} & \textbf{Description} \\
\midrule
\pyfile{seeds.py}       & \pyfile{ga.py}               & \texttt{list[str]}   & 15 initial prompt strings \\
\pyfile{dataset.py}     & \pyfile{ga.py}, \pyfile{main.py} & \texttt{list[dict]}  & Problem dicts (id, question, inputs, outputs, rating) \\
\pyfile{llm.py}         & \pyfile{ga.py}, \pyfile{main.py} & \texttt{str}         & Raw LLM response text \\
\pyfile{evaluator.py}   & \pyfile{ga.py}, \pyfile{main.py} & \texttt{str}         & Extracted C++ code string \\
\pyfile{evaluator.py}   & \pyfile{ga.py}, \pyfile{main.py} & \texttt{dict}        & Verbose test result (score, compile, test cases) \\
\pyfile{ga.py}          & \pyfile{main.py}             & \texttt{Individual}  & Best chromosome (prompt + fitness + metadata) \\
\pyfile{ga.py}          & \pyfile{main.py}             & \texttt{list[dict]}  & History: best/mean/worst/std/cv per generation \\
\pyfile{llm.py}         & \pyfile{cost\_tracker.py}   & \texttt{int, float}  & token counts + cost\_usd per API call \\
\pyfile{cost\_tracker.py} & \pyfile{main.py}           & \texttt{Path}        & Path to written cost JSON file \\
\pyfile{main.py}        & \pyfile{genealogy.py}        & \texttt{Path, str}   & run\_dir + best uid \\
\bottomrule
\end{tabular}
\caption{Data flowing between modules.}
\end{table}

\section{The API Call Budget}

One full run (Pop=5, Gen=12, Evo=65, Holdout=200) makes the following API calls:

\begin{table}[h]
\centering
\begin{tabular}{@{}lrrl@{}}
\toprule
\textbf{Call type} & \textbf{Count} & \textbf{Model} & \textbf{Formula} \\
\midrule
Fitness (gen 0 seeds)   & 325    & Flash & 5 seeds $\times$ 65 problems \\
Fitness (gens 1--12)    & 3{,}900 & Flash & 12 gens $\times$ 5 offspring $\times$ 65 problems \\
Holdout evaluation      & 200    & Flash & 1 best prompt $\times$ 200 problems \\
Crossover               & $\leq$48 & Pro & 12 gens $\times$ 5 offspring $\times$ 80\% \\
Mutate inject           & $\leq$9  & Pro & 12 gens $\times$ 3 mut-eligible $\times$ 25\% \\
Mutate delete           & $\leq$5  & Pro & 12 gens $\times$ 3 mut-eligible $\times$ 15\% \\
Mutate rephrase         & $\leq$4  & Pro & 12 gens $\times$ 3 mut-eligible $\times$ 10\% \\
\midrule
\textbf{Total Flash}    & $\approx$\textbf{4{,}425} & --- & --- \\
\textbf{Total Pro}      & $\approx$\textbf{66} & --- & --- \\
\bottomrule
\end{tabular}
\caption{Approximate API call counts per full run. ``Mut-eligible'' = offspring\_idx $\geq$ CROSSOVER\_ONLY\_SLOTS=2.}
\end{table}

\begin{conceptbox}{CROSSOVER\_ONLY\_SLOTS}
The first \fn{CROSSOVER\_ONLY\_SLOTS=2} offspring bred each generation receive crossover
only --- no mutation operators fire on them.  This preserves the best genetic material
more faithfully in the early slots while still allowing recombination.  The remaining 3
offspring per generation can receive any combination of inject, delete, and rephrase
mutations in addition to crossover.
\end{conceptbox}

\section{Fitness Score: What Is Being Optimised?}

The fitness of one prompt = the average partial-credit score across 65 Codeforces problems.
The partial credit for one problem = (test cases passed) / (total test cases).

\[
\text{fitness}(p) = \frac{1}{65} \sum_{i=1}^{65}
  \frac{\text{tests passed on problem } i \text{ using prompt } p}{\text{total tests on problem } i}
\]

A fitness of 1.0 means the prompt caused the LLM to write C++ code that passed \emph{every}
test case on \emph{every} problem.  Infrastructure failures (all retries exhausted) score $-1.0$
and are included in the mean, creating a small penalty for unreliable prompts.


% ==========================================================================
%  CHAPTER 3 --- config.py
% ==========================================================================
\chapter{Configuration: \pyfile{config.py}}

\section{Purpose}

Every hyperparameter and file path lives in this one file.  No other file hard-codes any
numbers or paths.  Authentication is done via a service account JSON file
(\fn{service-account.json}) that grants access to Google Cloud Vertex AI --- no interactive
login or API key string is required.

\section{File at a Glance}

\begin{table}[h]
\centering
\begin{tabular}{@{}llp{7cm}@{}}
\toprule
\textbf{Name} & \textbf{Type} & \textbf{Description} \\
\midrule
\fn{ROOT}                    & \texttt{Path}  & Project root (one level above \texttt{Code/}) \\
\fn{APPS\_TRAIN}             & \texttt{Path}  & Path to APPS/Codeforces training data \\
\fn{APPS\_TEST}              & \texttt{Path}  & Path to APPS/Codeforces test data \\
\fn{RESULTS\_DIR}            & \texttt{Path}  & Where run output folders are created \\
\fn{SERVICE\_ACCOUNT\_PATH}  & \texttt{Path}  & Path to service-account.json credential file \\
\fn{GCP\_REGION}             & \texttt{str}   & Vertex AI region (\texttt{"us-central1"}) \\
\fn{TARGET\_MODEL}           & \texttt{str}   & Gemini model for code generation (\texttt{gemini-2.5-flash}) \\
\fn{OPTIMIZER\_MODEL}        & \texttt{str}   & Gemini model for crossover/mutation (\texttt{gemini-2.5-pro}) \\
\fn{TARGET\_TEMPERATURE}     & \texttt{float} & Temperature for code generation (0.0, enforced) \\
\fn{OPTIMIZER\_TEMPERATURE}  & \texttt{float} & Temperature for operators (0.7, enforced) \\
\fn{API\_FALLBACK\_MAX\_TOKENS\_CODE} & \texttt{int} & Max output tokens for code gen (8192) \\
\fn{API\_FALLBACK\_MAX\_TOKENS\_OPS}  & \texttt{int} & Max output tokens for operators (700) \\
\fn{EVOLUTION\_SIZE}         & \texttt{int}   & Problems used for fitness during GA (65) \\
\fn{HOLDOUT\_SIZE}           & \texttt{int}   & Problems reserved for final test (200) \\
\fn{RANDOM\_SEED}            & \texttt{int}   & Seed for all random decisions (42) \\
\fn{EARLY\_STOP\_FAILURES}   & \texttt{int}   & Stop after N consecutive test failures (3) \\
\fn{POPULATION\_SIZE}        & \texttt{int}   & Chromosomes alive at any time (5) \\
\fn{GENERATIONS}             & \texttt{int}   & Breeding cycles (12) \\
\fn{ELITISM\_COUNT}          & \texttt{int}   & Top individuals copied unchanged (0) \\
\fn{CROSSOVER\_ONLY\_SLOTS}  & \texttt{int}   & First N offspring get crossover only (2) \\
\fn{TOURNAMENT\_SIZE}        & \texttt{int}   & Contestants per tournament selection (3) \\
\fn{CROSSOVER\_RATE}         & \texttt{float} & Probability crossover fires per offspring (0.8) \\
\fn{MUT\_INJECT\_PROB}       & \texttt{float} & Probability inject mutation fires (0.25) \\
\fn{MUT\_DELETE\_PROB}       & \texttt{float} & Probability delete mutation fires (0.15) \\
\fn{MUT\_REPHRASE\_PROB}     & \texttt{float} & Probability rephrase mutation fires (0.10) \\
\fn{CONVERGENCE\_CV\_THRESHOLD} & \texttt{float} & CV $<$ this $\Rightarrow$ stagnant (0.01) \\
\fn{CONVERGENCE\_PATIENCE}   & \texttt{int}   & Stagnant generations before early stop (3) \\
\fn{CONVERGENCE\_MIN\_MEAN}  & \texttt{float} & Don't check CV if mean fitness below this (0.01) \\
\fn{EXEC\_TIMEOUT}           & \texttt{int}   & Seconds per test-case binary run (2) \\
\fn{COMPILE\_TIMEOUT}        & \texttt{int}   & Seconds allowed for g++ compilation (15) \\
\fn{PARALLEL\_EVALS}         & \texttt{int}   & Problems evaluated simultaneously (5) \\
\bottomrule
\end{tabular}
\caption{\pyfile{config.py}: all exported names.}
\end{table}

\section{Full Source Code}

\begin{lstlisting}[style=pycode, caption={\pyfile{config.py} --- complete source.}]
"""
config.py -- All hyperparameters and paths in one place.
Edit here; everything else reads from this module.
"""

from pathlib import Path

# ---------------------------------------------------------------------------
# Paths
# ---------------------------------------------------------------------------
ROOT         = Path(__file__).parent
APPS_TRAIN   = ROOT / "Dataset" / "APPS" / "train"
APPS_TEST    = ROOT / "Dataset" / "APPS" / "test"
RESULTS_DIR  = Path(__file__).parent / "results"

# ---------------------------------------------------------------------------
# Dataset
# ---------------------------------------------------------------------------
# 305 pre-selected Codeforces problems (stdin/stdout, full C++ programs).
# Evolution pool: 65 stratified (60 from tiers 800-1700 + 5 harder 1800/1900).
EVOLUTION_SIZE  = 65    # fixed pool for fitness evaluation during GA
HOLDOUT_SIZE    = 200   # held-out for final within-distribution test
RANDOM_SEED     = 42
EARLY_STOP_FAILURES = 3   # stop running test cases after this many consecutive failures

# ---------------------------------------------------------------------------
# Genetic Algorithm
# ---------------------------------------------------------------------------
POPULATION_SIZE      = 5
GENERATIONS          = 12
ELITISM_COUNT        = 0   # no frozen elites -- all slots bred via crossover+mutation
CROSSOVER_ONLY_SLOTS = 2   # first N offspring get crossover only (no mutation)
TOURNAMENT_SIZE      = 3   # contestants per tournament selection draw

CROSSOVER_RATE      = 0.8

# Independent per-individual mutation probabilities (can all apply at once)
MUT_INJECT_PROB   = 0.25
MUT_DELETE_PROB   = 0.15
MUT_REPHRASE_PROB = 0.10

# ---------------------------------------------------------------------------
# Google Cloud Vertex AI -- Gemini backend
# ---------------------------------------------------------------------------
SERVICE_ACCOUNT_PATH = Path(__file__).parent / "service-account.json"
GCP_REGION           = "us-central1"

# Model being optimised -- generates code, evaluated for fitness.
TARGET_MODEL    = "gemini-2.5-flash"

# Model used for crossover / mutation meta-prompts.
OPTIMIZER_MODEL = "gemini-2.5-pro"

TARGET_TEMPERATURE    = 0.0
OPTIMIZER_TEMPERATURE = 0.7

API_FALLBACK_MAX_TOKENS_CODE = 8192   # max output tokens for code generation
API_FALLBACK_MAX_TOKENS_OPS  = 700    # crossover / mutation outputs are short

# ---------------------------------------------------------------------------
# Convergence detection  (CV = std / mean, per generation)
# ---------------------------------------------------------------------------
CONVERGENCE_CV_THRESHOLD = 0.01   # 1% variation -> converged
CONVERGENCE_PATIENCE     = 3      # consecutive stagnant generations to confirm
CONVERGENCE_MIN_MEAN     = 0.01   # only check CV when mean fitness exceeds this

# ---------------------------------------------------------------------------
# Code execution  (g++ compilation + binary runs)
# ---------------------------------------------------------------------------
EXEC_TIMEOUT    = 2    # seconds per test-case binary run
COMPILE_TIMEOUT = 15   # seconds allowed for g++ compilation
PARALLEL_EVALS  = 5    # problems evaluated simultaneously per individual
\end{lstlisting}

\section{Block-by-Block Explanation}

\subsection{Path Resolution}

\begin{lstlisting}[style=chain]
ROOT        = Path(__file__).parent
               -- __file__ = ".../Code/config.py"
               -- .parent  = ".../Code/"
               -- ROOT     = ".../Code/"

APPS_TRAIN  = ROOT / "Dataset" / "APPS" / "train"
            = ".../Code/Dataset/APPS/train"
APPS_TEST   = ROOT / "Dataset" / "APPS" / "test"
            = ".../Code/Dataset/APPS/test"
RESULTS_DIR = Path(__file__).parent / "results"
            = ".../Code/results"

SERVICE_ACCOUNT_PATH = Path(__file__).parent / "service-account.json"
            = ".../Code/service-account.json"
\end{lstlisting}

Both \fn{APPS\_TRAIN} and \fn{APPS\_TEST} are verified by \fn{main.main()} before the GA
starts.  \fn{SERVICE\_ACCOUNT\_PATH} is checked for existence and parsed for \fn{project\_id}
as the very first pre-flight step.

\subsection{Dataset Constants}

\begin{itemize}[nosep]
  \item \fn{EVOLUTION\_SIZE = 65} --- The GA evaluates every prompt on exactly these
    65 problems.  Keeping the set fixed across all generations is critical: if the
    problems changed, a prompt might score higher simply because it got easier problems.
    The actual set is determined by stratified sampling (\fn{dataset.split()}), not a
    simple first-65 slice.
  \item \fn{HOLDOUT\_SIZE = 200} --- After the GA finishes, the \emph{best} prompt is
    tested on these 200 \emph{never-seen} problems to measure generalisation.
  \item \fn{RANDOM\_SEED = 42} --- Passed to \fn{dataset.split()} and \fn{ga.run()}'s
    \texttt{random.Random} instance to make the experiment fully reproducible.
  \item \fn{EARLY\_STOP\_FAILURES = 3} --- In \fn{evaluator.run\_tests\_verbose()}: after
    3 consecutive test failures, the test loop exits early to save time.
\end{itemize}

\subsection{Genetic Algorithm Constants}

\begin{conceptbox}{No elitism (ELITISM\_COUNT = 0)}
Unlike classic GA designs, the current version does not freeze any individuals across
generations.  Every generation, all 5 slots are filled by new offspring from crossover
and mutation.  Elitism was removed because with a small population of 5, freezing 2
elites leaves only 3 slots for exploration --- too restrictive.  The best individual
is still preserved implicitly because tournament selection strongly favours it as a
parent.
\end{conceptbox}

\begin{conceptbox}{CROSSOVER\_ONLY\_SLOTS = 2}
The first 2 offspring bred each generation get crossover but no mutation operators.
This creates two ``conservative'' children that faithfully combine parent strategies.
The remaining 3 offspring also receive the three independent mutation operators, which
introduce novel strategies.  This hybrid balances exploitation (crossover-only) and
exploration (mutation).
\end{conceptbox}

\begin{conceptbox}{Tournament size: exploration vs.\ exploitation}
\fn{TOURNAMENT\_SIZE = 3} means we draw 3 random contestants and select the best.
With only 5 individuals in the population, a tournament of 3 draws a large fraction
of the pool, creating strong selection pressure toward the current best.
\end{conceptbox}

\subsection{Vertex AI Constants}

\fn{SERVICE\_ACCOUNT\_PATH} points to a Google Cloud service account JSON file that grants
access to the Vertex AI API.  The file is read once at startup by \fn{llm.\_get\_api\_client()};
credentials are then held in a singleton \fn{google.genai.Client} object shared across all threads.

\fn{GCP\_REGION = "us-central1"} is the Vertex AI region.  Both \fn{gemini-2.5-flash} and
\fn{gemini-2.5-pro} are confirmed available in this region.

Unlike the previous Claude CLI backend, temperature is now \emph{enforced} via
\fn{GenerateContentConfig(temperature=...)} in every API call.  \fn{TARGET\_TEMPERATURE = 0.0}
gives deterministic, reproducible code generation; \fn{OPTIMIZER\_TEMPERATURE = 0.7} adds
creative variation to crossover and mutation.

\subsection{Convergence Detection}

The GA can stop early if the population becomes statistically homogeneous:

\begin{itemize}[nosep]
  \item \fn{CONVERGENCE\_CV\_THRESHOLD = 0.01}: CV (coefficient of variation = std/mean)
    below 1\% means all individuals are performing nearly identically.
  \item \fn{CONVERGENCE\_PATIENCE = 3}: the CV must stay below the threshold for 3
    \emph{consecutive} generations before early stopping triggers.
  \item \fn{CONVERGENCE\_MIN\_MEAN = 0.01}: CV is only meaningful when mean fitness is
    positive.  If all individuals are failing (mean $\approx$ 0), the early-stop check
    is skipped to avoid false convergence signals.
\end{itemize}


% ==========================================================================
%  CHAPTER 4 --- seeds.py
% ==========================================================================
\chapter{Seed Prompts: \pyfile{seeds.py}}

\section{Purpose}

The GA must start with a population.  \pyfile{seeds.py} provides 15 hand-crafted
instruction prompts --- the ``generation zero'' chromosomes.  Each emphasises a
different cognitive strategy for solving C++ competitive programming problems, ensuring
the GA starts with genuinely diverse material to crossover and mutate.  The pool of 15
seeds is larger than the current \fn{POPULATION\_SIZE = 5}; \fn{main.py} uses only
\fn{SEED\_PROMPTS[:POPULATION\_SIZE]}, so the first 5 seeds form the initial population.

\begin{conceptbox}{Why diversity in seeds matters}
If all 5 seeds said the same thing, crossover would merely re-shuffle identical text.
By providing 15 different strategies --- edge-case thinking, algorithm identification,
pseudocode-first, constraint-awareness, etc.\ --- crossover can combine genuinely
different approaches from the start, providing the GA with rich exploratory material.
\end{conceptbox}

\section{File at a Glance}

\begin{table}[h]
\centering
\begin{tabular}{@{}llp{8cm}@{}}
\toprule
\textbf{Name} & \textbf{Type} & \textbf{Description} \\
\midrule
\fn{SEED\_PROMPTS} & \texttt{list[str]} & 15 diverse C++ instruction strings (gen-0 chromosomes) \\
\bottomrule
\end{tabular}
\caption{\pyfile{seeds.py}: exported names.}
\end{table}

\section{The 15 Strategies at a Glance}

\begin{table}[h]
\centering
\begin{tabular}{@{}rp{3.2cm}p{7.8cm}@{}}
\toprule
\textbf{\#} & \textbf{Strategy} & \textbf{Core instruction focus} \\
\midrule
1  & Structured analysis     & Analyse input/output format, constraints, edge cases before coding \\
2  & Terse and direct        & Full C++ solution; g++ -std=c++23; no commentary \\
3  & Competitive programmer  & Find algorithm pattern (greedy, DP, graph, binary search, math) \\
4  & Step-by-step reasoning  & 4-step: understand, identify, boundary conditions, implement \\
5  & Example-driven          & Trace sample inputs manually before writing code \\
6  & Conversational          & Full C++ source with main; compile and handle all edge cases \\
7  & Edge-case defender      & n=0, n=1, all equal, max values, integer overflow \\
8  & Algorithm-first         & Determine algorithm + complexity before writing code \\
9  & Specification-heavy     & 5-point spec: cin, cout, g++ -O2 -std=c++23, all cases, exact format \\
10 & Pseudocode-then-code    & English sketch first, then translate to C++ \\
11 & Constraint-aware        & If n can be $10^5$, O($n^2$) will TLE \\
12 & Defensive / robust      & Adversarial inputs at extreme constraint boundaries \\
13 & Readability-first       & Meaningful variable names; correct beats clever \\
14 & Verify-against-examples & Mental verification against sample inputs before finalising \\
15 & Direct and focused      & Entire C++ program; reads stdin; writes stdout; passes all tests \\
\bottomrule
\end{tabular}
\caption{The 15 seed strategies. Prompts 1--5 form the initial population (POPULATION\_SIZE=5).}
\end{table}

\section{The Safety Assert}

\begin{lstlisting}[style=example]
assert len(SEED_PROMPTS) >= 5, \
    f"Need at least 5 seed prompts; got {len(SEED_PROMPTS)}"
\end{lstlisting}

If you accidentally remove seeds below the minimum, Python immediately raises
\texttt{AssertionError} when \pyfile{seeds.py} is imported --- before any API calls are made.
The threshold is 5 (= \fn{POPULATION\_SIZE}) rather than exactly 15, to allow the seed
pool to be extended without changing the assert.

\subsection{How Seeds Become Individuals}

In \fn{ga.run()} [\pyfile{ga.py}]:
\begin{lstlisting}[style=example]
seeds = SEED_PROMPTS[:config.POPULATION_SIZE]   # first 5 of the 15 seeds
population = [Individual(prompt=p, operators="seed") for p in seeds]
# Each seed string becomes an Individual with:
#   prompt              = the seed string
#   fitness             = -1.0  (unevaluated sentinel)
#   uid                 = random 8-char hex (e.g. "a3f7c2b1")
#   parent_a            = ""  (empty: seeds have no parents)
#   parent_b            = ""
#   operators           = "seed"
#   intermediate_steps  = {}
#   failure_examples    = []
\end{lstlisting}


% ==========================================================================
%  CHAPTER 5 --- dataset.py
% ==========================================================================
\chapter{Dataset Loading: \pyfile{dataset.py}}

\section{Purpose}

This file loads 305 pre-selected Codeforces competitive programming problems from disk
and splits them into three non-overlapping groups: the \textbf{evolution pool} (65
problems for fitness evaluation during the GA), the \textbf{held-out set} (200 problems
for the final generalisation test), and the \textbf{rest} (remaining $\approx$40 problems).

The split is \textbf{stratified}: harder problems (rating $\geq$ 1800) are always included
in the evolution pool; regular tiers 800--1700 are drawn proportionally per
\fn{\_EVOLUTION\_PER\_TIER} counts.  This is fundamentally different from the old simple
shuffle.

\section{File at a Glance}

\begin{table}[h]
\centering
\begin{tabular}{@{}llp{7.5cm}@{}}
\toprule
\textbf{Item} & \textbf{Type} & \textbf{Description} \\
\midrule
\fn{\_SELECTED}              & \texttt{dict} & Hard-coded problem IDs by split then rating \\
\fn{\_EVOLUTION\_PER\_TIER}  & \texttt{dict} & How many problems to draw from each tier (800--1700) \\
\fn{\_RESERVED\_PER\_TIER}   & \texttt{dict} & How many leftover problems go into the reserved subfolder \\
\fn{load\_problems()}        & function & Load all 305 problems from disk; returns list of dicts \\
\fn{split(...)}              & function & Stratified split into evolution (65) / holdout (200) / rest \\
\fn{reserved\_train\_problems()} & function & Return 35 ``reserved'' train problems for the reserved set \\
\bottomrule
\end{tabular}
\caption{\pyfile{dataset.py}: all items.}
\end{table}

\section{Problem Distribution}

\begin{table}[h]
\centering
\begin{tabular}{@{}llrrrr@{}}
\toprule
\textbf{Tier} & \textbf{Rating} & \textbf{Train} & \textbf{Test} & \textbf{Total} & \textbf{Evo draw} \\
\midrule
Tier 1 & 800  & 29 & 21 & 50 & 7 \\
Tier 1 & 1000 & 10 & 40 & 50 & 6 \\
Tier 2 & 1100 & 14 & 19 & 33 & 6 \\
Tier 2 & 1200 & 23 & 10 & 33 & 7 \\
Tier 2 & 1300 & 23 & 11 & 34 & 7 \\
Tier 3 & 1400 & 27 &  5 & 32 & 7 \\
Tier 3 & 1500 & 33 &  0 & 33 & 7 \\
Tier 4 & 1600 & 17 &  0 & 17 & 6 \\
Tier 4 & 1700 & 18 &  0 & 18 & 7 \\
Tier 5 & 1800 &  3 &  0 &  3 & 3 (always) \\
Tier 5 & 1900 &  2 &  0 &  2 & 2 (always) \\
\midrule
\textbf{Total} & & 199 & 106 & 305 & \textbf{65} \\
\bottomrule
\end{tabular}
\caption{Problem distribution. Harder problems (1800/1900) are always in the evolution pool.}
\end{table}

\section{Function-by-Function Analysis}

\subsection{\fn{load\_problems()} $\to$ \fn{list[dict]}}

\begin{specbox}{load\_problems() -> list[dict]}
\textbf{Called by:} \fn{main.main()}\\
\textbf{Calls:} \fn{json.loads()}, \fn{Path.read\_text()}, \fn{Path.exists()}\\
\textbf{Returns:} \texttt{list[dict]} --- one dict per successfully loaded problem ($\leq$305)
\end{specbox}

Iterates over the hard-coded \fn{\_SELECTED} dict (split $\to$ rating $\to$ folder IDs).
For each folder it reads \fn{question.txt} and \fn{input\_output.json}.  Skips problems
with missing files, bad JSON, or empty test cases.

\begin{exbox}{A real problem dict}
\begin{lstlisting}[style=example]
{
    "id":       "train/0002",
    "question": "Given n integers on a line, compute their sum.\n"
                "Input: first line is n; second line has n integers.\n"
                "Output: a single integer.\nExample input:\n3\n1 2 3\n"
                "Example output:\n6",
    "inputs":   ["3\n1 2 3\n", "1\n0\n", "5\n-1 -2 3 4 5\n"],
    "outputs":  ["6\n",        "0\n",    "9\n"],
    "rating":   800,
}
# inputs[i]  = raw string piped to the compiled C++ binary's stdin
# outputs[i] = expected string on stdout (whitespace-normalised at compare time)
\end{lstlisting}
\end{exbox}

\subsection{\fn{split(problems, evolution\_size, holdout\_size, seed=42)}}

\begin{specbox}{split(...) -> (evolution[65], holdout[200], rest[~40])}
\textbf{Called by:} \fn{main.main()}\\
\textbf{Parameters:}\\
\hspace{1em}\param{problems} --- full list from \fn{load\_problems()} ($\leq$305 dicts)\\
\hspace{1em}\param{evolution\_size} --- accepted for API compat but actual size is stratified\\
\hspace{1em}\param{holdout\_size} --- size of the held-out test set (200)\\
\hspace{1em}\param{seed} --- RNG seed (42)\\
\textbf{Returns:} \texttt{tuple[list, list, list]} --- (evolution, holdout, rest)
\end{specbox}

The split is \textbf{stratified}, not a simple shuffle:

\begin{lstlisting}[style=example]
# Step 1: separate harder train problems (1800/1900) -- always in evolution
harder_train = [p for p in problems
                if p["id"].startswith("train/") and p["rating"] >= 1800]

# Step 2: group regular train problems by rating tier
regular_by_tier = defaultdict(list)
for p in problems:
    if p["id"].startswith("train/") and p["rating"] < 1800:
        regular_by_tier[p["rating"]].append(p)

# Step 3: draw per-tier using _EVOLUTION_PER_TIER counts (seed-42 shuffle)
rng = random.Random(seed)
evo_regular = []
for rating in sorted(regular_by_tier):
    tier_probs = regular_by_tier[rating].copy()
    rng.shuffle(tier_probs)                          # deterministic
    n = _EVOLUTION_PER_TIER.get(rating, 0)           # 6 or 7 depending on tier
    evo_regular.extend(tier_probs[:n])

evolution = harder_train + evo_regular               # 5 + 60 = 65 problems

# Step 4: everything else goes to holdout / rest
evo_ids   = {p["id"] for p in evolution}
remaining = [p for p in problems if p["id"] not in evo_ids]
rng.shuffle(remaining)
holdout   = remaining[:holdout_size]                 # 200 problems
rest      = remaining[holdout_size:]                 # ~40 problems
\end{lstlisting}

\begin{conceptbox}{Stratification vs.\ simple shuffle}
The old design shuffled all 300 problems uniformly and took the first 100.  This meant
harder problems (1700+ rating) might or might not appear in the evolution pool depending
on the random seed --- making the fitness signal unstable across experiments.  The new
stratified design guarantees:
\begin{itemize}[nosep]
  \item All 5 harder problems (1800/1900) are \emph{always} in the evolution pool.
  \item Each regular tier contributes a predictable count (6 or 7 problems).
  \item The fitness signal has a controlled, consistent difficulty distribution.
\end{itemize}
\end{conceptbox}

\subsection{\fn{reserved\_train\_problems(problems, seed=42)} $\to$ \fn{list[dict]}}

\begin{specbox}{reserved\_train\_problems(problems, seed=42) -> list[dict]}
\textbf{Called by:} \fn{setup\_reserved.py} (not the main GA loop)\\
\textbf{Returns:} 35 train problems (800--1700) for the reserved subfolder
\end{specbox}

Returns the problems that were \emph{not} selected for evolution but were still curated:
the 35 problems (from tiers 800--1700) drawn using \fn{\_RESERVED\_PER\_TIER} counts, taken
immediately after the evolution picks in the same tier shuffle order.  Together with the
65 active evolution problems this gives 100 curated train problems total; the remaining
205 are from the test split.


% ==========================================================================
%  CHAPTER 6 --- evaluator.py
% ==========================================================================
\chapter{Code Evaluation: \pyfile{evaluator.py}}

\section{Purpose}

When the LLM generates a response, we need to:
\begin{enumerate}[nosep]
  \item \textbf{Extract} the C++ code from the LLM's raw text response.
  \item \textbf{Compile} the extracted code with \fn{g++ -O2 -std=c++23}.
  \item \textbf{Execute} the compiled binary against each test case, piping the input
    string to stdin and comparing stdout against the expected output.
  \item \textbf{Score} the solution as (tests passed) / (total tests) with early stopping.
\end{enumerate}

This file has no LLM calls and no external dependencies beyond the Python standard library.
Process-group killing (\fn{\_kill\_group}) ensures that orphaned child processes from timeouts
never block the evaluator.

\section{File at a Glance}

\begin{table}[h]
\centering
\begin{tabular}{@{}llp{8cm}@{}}
\toprule
\textbf{Item} & \textbf{Type} & \textbf{Description} \\
\midrule
\fn{\_MAX\_IO\_LEN}           & constant & Max chars per test case in JSONL (500) \\
\fn{\_kill\_group(proc)}      & private  & Kill entire process group of proc (prevents orphans) \\
\fn{extract\_code(response)}  & function & Pull C++ code out of raw LLM text \\
\fn{\_normalise(text)}        & private  & Strip trailing whitespace per line \\
\fn{\_outputs\_match(a, e)}   & private  & Compare normalised actual vs.\ expected \\
\fn{\_compile(code)}          & private  & Write to temp .cpp; run g++ -std=c++23; return binary path \\
\fn{\_run\_one(bin, stdin, t)} & private & Run compiled binary with one input string \\
\fn{run\_tests\_verbose(...)}  & function & Full pipeline: compile + run + early-stop + detailed dict \\
\fn{run\_tests(...)}           & function & Thin wrapper returning only the float score \\
\bottomrule
\end{tabular}
\caption{\pyfile{evaluator.py}: all items.}
\end{table}

\section{Function-by-Function Analysis}

\subsection{\fn{\_kill\_group(proc)}}

\begin{specbox}{\_kill\_group(proc: subprocess.Popen) -> None}
\textbf{Called by:} \fn{\_compile()}, \fn{\_run\_one()} on timeout
\end{specbox}

Sends \fn{SIGKILL} to the entire process group of \fn{proc}:
\begin{lstlisting}[style=example]
def _kill_group(proc):
    try:
        os.killpg(os.getpgid(proc.pid), signal.SIGKILL)
    except (ProcessLookupError, PermissionError):
        try:
            proc.kill()   # fallback: kill only the direct process
        except Exception:
            pass
\end{lstlisting}

Without this, a timed-out C++ program that spawned child processes could leave them
running and writing to the pipe, causing the parent's \fn{communicate()} to hang.
Both \fn{\_compile()} and \fn{\_run\_one()} use \fn{preexec\_fn=os.setsid} when launching
subprocesses, which places them in a new session (= their own process group), making
\fn{os.killpg} safe to call.

\subsection{\fn{extract\_code(response: str)} $\to$ \fn{str}}

Tries three strategies in priority order:
\begin{lstlisting}[style=example]
# Strategy 1: tagged C++ fence
m = re.search(r"```(?:cpp|c\+\+)\s*\n([\s\S]*?)```", response, re.IGNORECASE)
if m: return m.group(1).strip()

# Strategy 2: any fence
m = re.search(r"```\s*\n?([\s\S]*?)```", response)
if m: return m.group(1).strip()

# Strategy 3: raw text (assume entire reply is code)
return response.strip()
\end{lstlisting}

Because both Gemini models are configured with a \fn{\_CODE\_GEN\_DIRECTIVE} that
explicitly forbids markdown fences, the LLM typically returns raw C++ code (strategy 3).
Strategies 1 and 2 handle the cases where the model ignores the directive.

\subsection{\fn{\_compile(code: str)} $\to$ \fn{(bin\_path | None, error\_msg)}}

\begin{specbox}{\_compile(code) -> (str | None, str)}
\textbf{Calls:} \fn{tempfile.mkstemp()}, \fn{subprocess.Popen(["g++", "-O2", "-std=c++23", ...])}
\end{specbox}

\begin{lstlisting}[style=chain]
_compile(code):
  src_fd, src_path = tempfile.mkstemp(suffix=".cpp")
  bin_path = src_path[:-4]   -- strip .cpp extension

  write code to src_path via os.fdopen(src_fd)

  proc = subprocess.Popen(
      ["g++", "-O2", "-std=c++23", "-o", bin_path, src_path],
      stdout=PIPE, stderr=PIPE, text=True,
      preexec_fn=os.setsid)              -- new session -> own process group

  try:
      _, stderr = proc.communicate(timeout=COMPILE_TIMEOUT)  -- 15s limit
  except subprocess.TimeoutExpired:
      _kill_group(proc); proc.communicate()
      return None, "compilation timed out"

  if proc.returncode != 0:
      return None, stderr[:500]          -- compile error text

  return bin_path, ""                    -- success: binary exists

  finally:
      os.unlink(src_path)                -- always delete .cpp
      if not _success: os.unlink(bin_path)  -- delete binary only on failure
\end{lstlisting}

\begin{warnbox}{g++ -std=c++23}
The compiler standard was upgraded from \fn{c++17} to \fn{c++23} to support modern
Codeforces solutions that use C++23 features (\fn{std::print}, \fn{std::ranges}, etc.).
The seed prompts also reference \fn{c++23} in several places (seeds 2 and 9).
g++ 13+ (available on Ubuntu 24.04) supports \fn{-std=c++23}.
\end{warnbox}

\subsection{\fn{\_run\_one(bin\_path, stdin\_str, timeout)} $\to$ \fn{Optional[str]}}

\begin{lstlisting}[style=chain]
_run_one(bin_path, stdin_str, timeout):
  proc = subprocess.Popen(
      [bin_path],
      stdin=PIPE, stdout=PIPE, stderr=DEVNULL,
      text=True,
      preexec_fn=os.setsid)          -- own process group for clean kill

  try:
      stdout, _ = proc.communicate(input=stdin_str, timeout=timeout)
      return stdout
  except subprocess.TimeoutExpired:
      _kill_group(proc); proc.communicate()
      return None                    -- timeout -> test case fails
  except Exception:
      return None
\end{lstlisting}

\subsection{\fn{run\_tests\_verbose(code, inputs, outputs, timeout=2)} $\to$ \fn{dict}}

\begin{specbox}{run\_tests\_verbose(...) -> dict}
\textbf{Called by:} \fn{ga.\_eval\_code\_for\_problem()}; \fn{main.\_holdout\_one()}\\
\textbf{Calls:} \fn{\_compile()}, \fn{\_run\_one()}, \fn{\_outputs\_match()}
\end{specbox}

Returns a rich dict suitable for JSONL logging:
\begin{lstlisting}[style=chain]
{
  "score":          float,        -- passed / total
  "reason":         str,          -- "empty_code" | "compile_fail" | "no_tests"
                                  --   | "perfect" | "partial_pass"
                                  --   | "all_fail" | "all_fail_early_stop"
  "compile":        {"success": bool, "error": str|None, "time_s": float},
  "test_cases":     [{"index": int, "passed": bool, "actual": str|None,
                       "expected": str, "time_s": float, "timeout": bool}, ...],
  "tests_total":    int,          -- total available test cases
  "tests_run":      int,          -- how many actually executed (early-stop)
  "early_stopped":  bool,
  "total_test_time_s": float,
}
\end{lstlisting}

\begin{conceptbox}{Early stopping}
After \fn{EARLY\_STOP\_FAILURES=3} consecutive failures, the test loop exits.
The score denominator is still \fn{len(inputs)} (total), so skipped tests count as
failures.  This saves time on clearly broken solutions without giving them a free pass.
\end{conceptbox}

\fn{run\_tests()} is a thin backward-compat wrapper:
\begin{lstlisting}[style=example]
def run_tests(code, inputs, outputs, timeout=2) -> float:
    return run_tests_verbose(code, inputs, outputs, timeout)["score"]
\end{lstlisting}


% ==========================================================================
%  CHAPTER 7 --- cost_tracker.py
% ==========================================================================
\chapter{Cost Tracking: \pyfile{cost\_tracker.py}}

\section{Purpose}

Logs token counts and estimated costs for every Vertex AI API call.  The tracker
\emph{streams} each call record to a JSONL file immediately (no in-memory list), keeping
memory usage flat regardless of run length.  Running totals are maintained as simple
scalars.  A summary JSON is written to \fn{Costs/} at the end of each run and also
live-updated after every call for real-time monitoring.

\begin{conceptbox}{Real charges}
The project calls Google Cloud Vertex AI, which bills real charges to your GCP account.
The \fn{PRICING} table lists the actual Gemini rates (USD per million tokens) as of
April 2026.  The \$300 GCP free credit applies to these charges.  Thinking tokens
(internal chain-of-thought) are billed at the same output token rate as regular output
tokens and are aggregated into \fn{output\_tokens} before cost calculation.
\end{conceptbox}

\section{The PRICING Table}

\begin{lstlisting}[style=example]
PRICING = {
    # model: (input_$/MTok, output_$/MTok)
    # Gemini on Vertex AI (active models)
    "gemini-2.5-flash":     (0.15,  0.60),
    "gemini-2.5-pro":       (1.25, 10.00),
    "gemini-2.0-flash":     (0.10,  0.40),
    "gemini-2.0-flash-001": (0.10,  0.40),
    # Claude on Vertex AI (kept for reference)
    "claude-haiku-4-5-20251001": (1.00,  5.00),
    "claude-sonnet-4-6":         (3.00, 15.00),
}
\end{lstlisting}

In practice only \fn{gemini-2.5-flash} and \fn{gemini-2.5-pro} are used.

\section{\fn{CallRecord} --- One Call's Data}

\begin{lstlisting}[style=example]
@dataclass
class CallRecord:
    phase: str          # "fitness", "crossover", "mutate_inject", etc.
    model: str
    input_tokens: int
    output_tokens: int  # already includes thinking tokens (summed by llm.py)
    cost_usd: float | None = None   # set directly from API response

    @property
    def total_cost(self) -> float:
        if self.cost_usd is not None:
            return self.cost_usd   # prefer the value passed in from llm.py
        rate_in, rate_out = PRICING.get(self.model, (0, 0))
        return (self.input_tokens * rate_in
              + self.output_tokens * rate_out) / 1_000_000
\end{lstlisting}

\section{\fn{CostTracker} --- The Streaming Singleton}

\begin{specbox}{class CostTracker}
Module-level singleton: \fn{tracker = CostTracker()}\\
Imported by \pyfile{llm.py} as: \fn{from cost\_tracker import tracker}
\end{specbox}

\begin{lstlisting}[style=chain]
CostTracker.set_log_path(path)
  Called by: main.main() before GA starts
  Action: sets path for streaming JSONL writes (calls.jsonl)

CostTracker.set_live_path(path)
  Called by: main.main()
  Action: path for live JSON updates inside the run directory

CostTracker.set_costs_live_path(path)
  Called by: main.main()
  Action: second live path (in Costs/ folder) -- updated after every call

CostTracker.restore_from_jsonl(path) -> int
  Called by: main.main() on resume
  Action: reloads running totals from an existing calls.jsonl
  Returns: number of lines loaded (for logging)

CostTracker.record(phase, model, input_tokens, output_tokens, cost_usd)
  Called by: llm._call_realtime_api() after every successful API response
  Thread-safe (internal _lock)
  Action:
    - updates running totals (_total_calls, _total_cost_usd, etc.)
    - writes one JSON line immediately to calls.jsonl:
        {"t":"HH:MM:SS", "phase":..., "model":..., "in":..., "out":...,
         "cost":..., "total_calls":..., "running_cost":...}
    - overwrites live summary files (costs_live.json, Costs/cost_<run>.json)

CostTracker.flush(extra=None) -> Path
  Called by: main.main() at the very end
  Action:
    - assembles summary from running totals (_phase_totals, _model_totals)
    - writes Code/Costs/cost_TIMESTAMP.json
    - returns Path of the written file
\end{lstlisting}

\begin{conceptbox}{Streaming vs.\ in-memory}
Each \fn{CallRecord} is written to disk immediately and the object is discarded.
Running totals (\fn{\_total\_calls}, \fn{\_total\_input\_tokens}, \fn{\_phase\_totals},
\fn{\_model\_totals}) are simple scalars/dicts updated in-place.  A run with 4,425 API
calls uses negligible memory for cost tracking --- there is no growing list of records.
\end{conceptbox}


% ==========================================================================
%  CHAPTER 8 --- llm.py
% ==========================================================================
\chapter{LLM Calls: \pyfile{llm.py}}

\section{Purpose}

The sole interface to Gemini models via Google Cloud Vertex AI.  All public functions
route through the private \fn{\_call\_realtime\_api()} wrapper, which:
\begin{enumerate}[nosep]
  \item Obtains a shared \fn{google.genai.Client} (lazy singleton, thread-safe).
  \item Calls \fn{client.models.generate\_content()} with the appropriate config.
  \item Tracks thinking tokens separately and combines them into the output token count.
  \item Retries on rate limits, server errors, and timeouts with exponential backoff.
  \item Records token counts and cost to the \fn{CostTracker} singleton.
\end{enumerate}

\section{File at a Glance}

\begin{table}[h]
\centering
\begin{tabular}{@{}llp{8.5cm}@{}}
\toprule
\textbf{Item} & \textbf{Type} & \textbf{Description} \\
\midrule
\fn{LLMCallError}          & exception & Raised when all retries fail for generate\_code() \\
\fn{set\_error\_log\_path()} & function & Set path for llm\_errors.log \\
\fn{\_get\_api\_client()}  & private  & Thread-safe lazy singleton: returns Vertex AI client \\
\fn{\_call\_realtime\_api()} & private & Actual API call with retry \& cost tracking \\
\fn{\_call()}              & private  & Router to \_call\_realtime\_api() \\
\fn{generate\_code()}      & function & Ask TARGET\_MODEL to solve a Codeforces problem \\
\fn{run\_fitness\_batch()}  & function & Stub: batch API not used (raises NotImplementedError) \\
\fn{crossover()}           & function & Ask OPTIMIZER\_MODEL to merge two prompts \\
\fn{mutate\_inject()}      & function & Ask OPTIMIZER to add one new strategy to a prompt \\
\fn{mutate\_delete()}      & function & Ask OPTIMIZER to remove the weakest sentence \\
\fn{mutate\_rephrase()}    & function & Ask OPTIMIZER to rephrase without changing meaning \\
\bottomrule
\end{tabular}
\caption{\pyfile{llm.py}: all items.}
\end{table}

\section{Output Directives}

Two constant strings are appended to every prompt at call time to enforce output format:

\begin{lstlisting}[style=example]
_CODE_GEN_DIRECTIVE = (
    "\n\nABSOLUTE OUTPUT REQUIREMENT -- NO EXCEPTIONS:\n"
    "Your entire response must consist of ONLY the raw C++ source code -- "
    "one complete file, nothing else. Use C++23.\n"
    "FORBIDDEN:\n"
    "  * Backtick code fences of any form (no ```, no `, no 'cpp' tags)\n"
    "  * Markdown formatting of any kind\n"
    "  * Any text before or after the code\n"
    "The FIRST character of your response must be the first character of the "
    "C++ source code. The LAST character must be the last character of the code."
)

_OPERATOR_DIRECTIVE = (
    "\n\nABSOLUTE OUTPUT REQUIREMENT -- NO EXCEPTIONS:\n"
    "Your entire response must consist of ONLY the plain text of the instruction "
    "prompt -- nothing else.\n"
    "FORBIDDEN: backtick fences, markdown, labels, meta-commentary."
)
\end{lstlisting}

These are appended at call time, keeping the chromosome text clean --- evolved prompts
never contain operational directives.

\section{\fn{\_get\_api\_client()}}

\begin{specbox}{\_get\_api\_client() -> google.genai.Client}
\textbf{Called by:} \fn{\_call\_realtime\_api()} (once per process, then cached)\\
\textbf{Thread-safe:} double-checked locking via \fn{\_api\_client\_lock}
\end{specbox}

\begin{lstlisting}[style=example]
def _get_api_client():
    global _api_client
    if _api_client is not None:
        return _api_client        # fast path: already initialised
    with _api_client_lock:
        if _api_client is not None:
            return _api_client    # double-check: another thread may have beaten us

        sa_info    = json.loads(config.SERVICE_ACCOUNT_PATH.read_text())
        project_id = sa_info["project_id"]
        credentials = service_account.Credentials.from_service_account_info(
            sa_info,
            scopes=["https://www.googleapis.com/auth/cloud-platform"],
        )
        _api_client = genai.Client(
            vertexai=True,
            project=project_id,
            location=config.GCP_REGION,        # "us-central1"
            credentials=credentials,
            http_options=HttpOptions(timeout=120_000),  # 120s per request
        )
    return _api_client
\end{lstlisting}

The 120-second HTTP timeout is set at the client level.  If any individual API request
takes longer, the SDK raises a timeout exception, which the retry loop catches and
handles with exponential backoff.

\section{\fn{\_call\_realtime\_api()} --- Core API Wrapper}

\begin{specbox}{\_call\_realtime\_api(phase, model, system, user, max\_tokens, retries=10) -> Optional[str]}
\textbf{Returns:} response text on success; \fn{""} if all retries return empty; \fn{None} if all retries fail
\end{specbox}

\begin{lstlisting}[style=chain]
_call_realtime_api(phase, model, system, user, max_tokens, retries=10):

  client = _get_api_client()

  for attempt in range(10):
    t0 = time.monotonic()
    try:
      is_fitness = (phase == "fitness")
      response = client.models.generate_content(
          model=model,
          contents=user,
          config=GenerateContentConfig(
              system_instruction=system,
              max_output_tokens=max_tokens,
              temperature=0.0 if is_fitness else 0.7,
              thinking_config=ThinkingConfig(
                  thinking_budget=512 if is_fitness else 2048,
              ),
          ),
      )

      text = response.text or ""

      # --- token accounting ---
      usage       = response.usage_metadata
      input_toks  = usage.prompt_token_count      or 0
      output_toks = usage.candidates_token_count  or 0
      think_toks  = usage.thoughts_token_count    or 0
      output_toks += think_toks    # thinking billed at output rate; merge

      cost = _compute_gemini_cost(model, input_toks, output_toks)
      tracker.record(phase=phase, model=model,
                     input_tokens=input_toks, output_tokens=output_toks,
                     cost_usd=cost)

      if not text.strip():          -- empty response
          wait = 5 * (2**attempt) + random.uniform(0, 3)
          time.sleep(wait); continue

      return text                   -- SUCCESS

    except Exception as exc:
      -- classify: is_rate (429/RESOURCE_EXHAUSTED), is_server (500/503),
      --           is_timeout ("timeout" / "timed out")
      -- all retryable: wait = 5 * 2^attempt + jitter; time.sleep(wait)
      -- non-retryable (auth, 404, etc.): still retry with same backoff

  return None                       -- all retries exhausted
\end{lstlisting}

\begin{conceptbox}{Thinking token accounting}
Gemini 2.5 models support extended thinking.  The model's internal chain-of-thought
appears in \fn{usage\_metadata.thoughts\_token\_count} but is \emph{not} included in
\fn{candidates\_token\_count}.  The code explicitly adds them:
\fn{output\_toks += think\_toks} before computing cost.  This means the cost figure
includes both the visible response and the hidden reasoning.
\end{conceptbox}

\section{\fn{generate\_code(system\_prompt, question)} $\to$ \fn{str}}

\begin{specbox}{generate\_code(system\_prompt, question) -> str}
\textbf{Called by:} \fn{ga.\_eval\_one\_problem()}; \fn{main.\_holdout\_one()}\\
\textbf{Raises:} \fn{LLMCallError} if all retries fail (returns \fn{None} from \fn{\_call})
\end{specbox}

\begin{lstlisting}[style=example]
def generate_code(system_prompt, question) -> str:
    user_msg = _build_fitness_user_msg(question)
    # _build_fitness_user_msg appends:
    #   "Write a complete, compilable C++ program that reads from stdin and
    #    writes the answer to stdout. Include main(). Compile target: g++ -O2
    #    -std=c++23. DO NOT write any explanation ... Your response must begin
    #    with #include or int main -- nothing before it."
    result = _call(
        phase="fitness",
        model=config.TARGET_MODEL,              # "gemini-2.5-flash"
        system=system_prompt + _CODE_GEN_DIRECTIVE,
        user=user_msg,
        api_max_tokens=config.API_FALLBACK_MAX_TOKENS_CODE,  # 8192
    )
    if result is None:
        raise LLMCallError(
            f"generate_code: all retries failed (model={config.TARGET_MODEL})"
        )
    return result
\end{lstlisting}

\section{\fn{crossover()} and the Three Mutation Operators}

\begin{lstlisting}[style=chain]
crossover(prompt_a, prompt_b) -> str      [phase="crossover", OPTIMIZER_MODEL]
  System: "You merge coding-instruction prompts. Output ONLY the merged prompt."
  User:   "Combine these two into one coherent prompt.
           Preserve the strongest strategies from each parent.
           Do not simply concatenate.
           --- Parent A ---  {parent_a}
           --- Parent B ---  {parent_b}"

mutate_inject(prompt, failure_examples=None) -> str
                                          [phase="mutate_inject", OPTIMIZER_MODEL]
  If failure_examples given:
      Appends: "Recent failure examples (address at least one):
               - Problem X (rating R): reason=compile_fail, ..."
  User: "Add one new, specific coding strategy ...
         Examples: 'validate your output', 'consider time complexity', ..."

mutate_delete(prompt) -> str              [phase="mutate_delete", OPTIMIZER_MODEL]
  Edge case FIRST:
      sentences = re.split(r"(?<=[.!?])\s+", prompt)
      if len(sentences) <= 1: return prompt   -- skip CLI call for single sentence
  User: "Remove one sentence -- the one that contributes the least."

mutate_rephrase(prompt) -> str            [phase="mutate_rephrase", OPTIMIZER_MODEL]
  User: "Rephrase using different words, keep same meaning and strategies."
\end{lstlisting}

\begin{conceptbox}{Failure examples in mutate\_inject}
When \fn{ga.breed()} calls \fn{mutate\_inject()}, it passes up to 5 failure examples
collected from the two parents during their fitness evaluation.  These describe actual
failures: compile errors, wrong answers, timeouts.  The Optimizer model sees this
context and is instructed to add a strategy targeting those specific failures.  This
creates a feedback loop: the GA learns from where the current best prompts fail.
\end{conceptbox}


% ==========================================================================
%  CHAPTER 9 --- ga.py
% ==========================================================================
\chapter{Genetic Algorithm: \pyfile{ga.py}}

\section{Purpose}

Implements the full evolutionary loop: seed evaluation, tournament selection,
crossover/mutation breeding, parallel fitness evaluation, convergence detection,
and crash-recovery checkpoints.  Also provides RNG state serialization for
deterministic resume.

\section{File at a Glance}

\begin{table}[h]
\centering
\begin{tabular}{@{}llp{8.5cm}@{}}
\toprule
\textbf{Item} & \textbf{Type} & \textbf{Description} \\
\midrule
\fn{Individual}              & dataclass & One chromosome (prompt + metadata) \\
\fn{rng\_state\_to\_dict()}  & function  & Serialize MT19937 RNG state to JSON-safe dict \\
\fn{rng\_state\_from\_dict()} & function & Deserialize RNG state dict back to tuple \\
\fn{\_eval\_code\_for\_problem()} & private & Extract + compile + test one response; write JSONL \\
\fn{\_eval\_one\_problem()}  & private   & Call Gemini for one problem; delegate to \_eval\_code\_for\_problem \\
\fn{compute\_fitness()}      & function  & Evaluate one individual on all 65 problems (parallel) \\
\fn{\_assign\_failure\_examples()} & private & Keep diverse subset of failure examples \\
\fn{tournament\_select()}    & function  & Pick best of TOURNAMENT\_SIZE random draws \\
\fn{roulette\_select()}      & function  & Fitness-proportionate (roulette-wheel) selection \\
\fn{breed()}                 & function  & Produce one offspring via crossover + mutation \\
\fn{run()}                   & function  & Main GA loop (init, evolve, converge, return best); accepts selection scheme \\
\bottomrule
\end{tabular}
\caption{\pyfile{ga.py}: all items.}
\end{table}

\section{\fn{Individual} --- The Chromosome Dataclass}

\begin{lstlisting}[style=example]
@dataclass
class Individual:
    prompt:             str             # the system prompt (the chromosome)
    fitness:            float = -1.0   # -1.0 = not yet evaluated
    uid:                str = field(default_factory=lambda: uuid.uuid4().hex[:8])
    parent_a:           str = ""       # uid of first parent (empty for seeds)
    parent_b:           str = ""       # uid of second parent (empty for seeds)
    operators:          str = ""       # e.g. "crossover+inject" or "seed"
    intermediate_steps: dict = field(default_factory=dict)
    # Records the prompt after each operator:
    #   {"after_crossover": "...", "after_inject": "...", "after_delete": "..."}
    failure_examples:   list = field(default_factory=list)
    # Up to 5 dicts describing problems where this individual scored < 1.0:
    #   [{"problem_id": "train/0002", "rating": 800, "reason": "compile_fail",
    #     "compile_error": "...", "expected": "...", "actual": "..."}, ...]
\end{lstlisting}

\section{\fn{\_eval\_code\_for\_problem()} and \fn{\_eval\_one\_problem()}}

\begin{specbox}{\_eval\_one\_problem(gen, individual, prob, idx, total) -> (float, dict|None)}
\textbf{Called by:} \fn{compute\_fitness()} via ThreadPoolExecutor\\
\textbf{Calls:} \fn{llm.generate\_code()}, then \fn{\_eval\_code\_for\_problem()}
\end{specbox}

\begin{lstlisting}[style=chain]
_eval_one_problem(gen, individual, prob, idx, total):
  time.sleep(random.uniform(0, 1))   -- stagger concurrent API requests

  t_llm = time.monotonic()
  try:
      response = llm.generate_code(individual.prompt, prob["question"])
  except llm.LLMCallError as exc:
      llm_time = monotonic() - t_llm
      return _eval_code_for_problem(gen, individual, prob, idx, total,
                                    None, llm_time)   -- None = infra failure

  llm_time = round(monotonic() - t_llm, 3)
  return _eval_code_for_problem(gen, individual, prob, idx, total,
                                response, llm_time)

_eval_code_for_problem(gen, individual, prob, idx, total, response_text, llm_time_s):
  if response_text is None:           -- infra failure path
      _log_evaluation({..., score:-1.0, reason:"infra_failure"})
      return -1.0, None

  code = evaluator.extract_code(response_text); del response_text

  if not code.strip():
      test_result = {score:0.0, reason:"empty_response", ...}
  else:
      test_result = evaluator.run_tests_verbose(code, prob["inputs"],
                                                prob["outputs"], EXEC_TIMEOUT)

  _log_evaluation({ts, gen, uid, operators, system_prompt, problem_id,
                   rating, llm_time, extracted_code, compile, test_cases,
                   tests_total, tests_run, early_stopped, score, reason})
  del entry, code

  -- build failure_example if score < 1.0 (not infra_failure):
  failure_example = {problem_id, rating, reason,
                     expected[:200], actual[:200], compile_error[:200]}

  return score, failure_example   -- (float, dict|None)
\end{lstlisting}

\section{\fn{compute\_fitness(gen, individual, problems)} $\to$ \fn{float}}

\begin{specbox}{compute\_fitness(...) -> float in [-1.0, 1.0]}
\textbf{Called by:} \fn{run()} --- 5 times at gen 0; 5 times per generation (gens 1--12)\\
\textbf{Parallelism:} ThreadPoolExecutor(max\_workers=PARALLEL\_EVALS=5)
\end{specbox}

\begin{lstlisting}[style=example]
def compute_fitness(gen, individual, problems):
    scores = [0.0] * len(problems)
    collected_failures = []
    _failures_lock = threading.Lock()

    with ThreadPoolExecutor(max_workers=config.PARALLEL_EVALS) as executor:
        futures = {
            executor.submit(_eval_one_problem, gen, individual, prob, i,
                            len(problems)): i
            for i, prob in enumerate(problems)
        }
        completed = 0
        for future in as_completed(futures):
            idx = futures[future]
            try:
                score, failure_example = future.result()
                scores[idx] = score
                if failure_example is not None:
                    with _failures_lock:
                        collected_failures.append(failure_example)
            except Exception as exc:
                scores[idx] = -1.0
            completed += 1
            if completed % 10 == 0 or completed == len(problems):
                _log(f"ind={individual.uid}  progress {completed}/{len(problems)}")

    _assign_failure_examples(individual, collected_failures)
    fitness = sum(scores) / len(scores)
    return fitness
\end{lstlisting}

\section{\fn{\_assign\_failure\_examples(individual, collected\_failures)}}

Keeps $\leq$5 failure examples, prioritising diversity of failure reason:

\begin{lstlisting}[style=example]
_reason_priority = {
    "compile_fail":        0,   # highest priority -- most informative
    "all_fail":            1,
    "all_fail_early_stop": 1,
    "partial_pass":        2,
}
# Sort by priority; then include an example of each reason type.
# At most 5 examples total; at most 3 per reason (to ensure diversity).
\end{lstlisting}

These failure examples are later passed to \fn{llm.mutate\_inject()} to focus new
strategy additions on the actual failure modes.

\section{RNG State Serialization}

\begin{lstlisting}[style=example]
def rng_state_to_dict(state: tuple) -> dict:
    version, internalstate, gauss_next = state
    return {
        "version":       version,
        "internalstate": list(internalstate),  # tuple -> list for JSON
        "gauss_next":    gauss_next,
    }

def rng_state_from_dict(d: dict) -> tuple:
    return (d["version"], tuple(d["internalstate"]), d["gauss_next"])
\end{lstlisting}

The MT19937 RNG state is saved in \fn{generation\_NN.json} checkpoints and restored on
resume, making tournament selection decisions deterministic across restarts.

\section{\fn{tournament\_select}, \fn{roulette\_select}, and \fn{breed}}

\begin{lstlisting}[style=chain]
tournament_select(population, rng):
  contestants = rng.sample(population, TOURNAMENT_SIZE=3)  -- no replacement
  return max(contestants, key=lambda ind: ind.fitness)

roulette_select(population, rng):
  fitnesses = [max(ind.fitness, 0.0) for ind in population]
  total = sum(fitnesses)
  if total == 0:
      return rng.choice(population)    -- fallback: uniform when all scores 0
  r = rng.uniform(0, total)
  cumulative = 0.0
  for ind, f in zip(population, fitnesses):
      cumulative += f
      if cumulative >= r:
          return ind
  return population[-1]                -- floating-point safety fallback

breed(parent_a, parent_b, rng, allow_mutation=True):

  [Collect combined failure examples from both parents (up to 5 unique)]

  [CROSSOVER]
  if rng.random() < 0.80:
      result = llm.crossover(parent_a.prompt, parent_b.prompt)
      prompt = result if result.strip() else rng.choice([a.prompt, b.prompt])
      ops = ["crossover"]
  else:
      prompt = rng.choice([parent_a.prompt, parent_b.prompt])
      ops = ["copy"]

  steps["after_crossover"] = prompt

  [MUTATION -- only when allow_mutation=True]
  if allow_mutation:
      if rng.random() < 0.25:
          result = llm.mutate_inject(prompt, failure_examples=combined)
          prompt = result if result.strip() else prompt
          ops.append("inject")
          steps["after_inject"] = prompt

      if rng.random() < 0.15:
          result = llm.mutate_delete(prompt)
          prompt = result if result.strip() else prompt
          ops.append("delete")
          steps["after_delete"] = prompt

      if rng.random() < 0.10:
          result = llm.mutate_rephrase(prompt)
          prompt = result if result.strip() else prompt
          ops.append("rephrase")
          steps["after_rephrase"] = prompt

  return Individual(prompt, parent_a=a.uid, parent_b=b.uid,
                    operators="+".join(ops), intermediate_steps=steps)
\end{lstlisting}

\section{\fn{run()} --- The Main GA Loop}

\begin{specbox}{run(seed\_prompts, problems, run\_dir, on\_generation, resume\_*, selection\_scheme, elitism\_count) -> (Individual, list[dict])}
\textbf{Called by:} \fn{main.main()}\\
\textbf{New parameters:}\\
\hspace{1em}\param{selection\_scheme} --- \fn{"tournament"} (default), \fn{"roulette"}, or \fn{"elitism"}\\
\hspace{1em}\param{elitism\_count} --- overrides \fn{config.ELITISM\_COUNT} when scheme is \fn{"elitism"} (pass 1)
\end{specbox}

\begin{lstlisting}[style=chain]
run(seed_prompts, problems, run_dir, on_generation,
    resume_population=None, resume_start_gen=0,
    resume_history=None, resume_rng_state=None,
    selection_scheme="tournament", elitism_count=None):

  _progress_log = run_dir / "progress.log"
  _eval_log_dir = run_dir

  -- resolve selection function
  _scheme = selection_scheme.lower()
  if _scheme == "roulette":
      _select = roulette_select
  else:                                -- "tournament" or "elitism"
      _select = tournament_select

  _elitism = elitism_count if elitism_count is not None else config.ELITISM_COUNT

  rng = random.Random(RANDOM_SEED=42)
  history = list(resume_history) if resume_history else []

  [INIT OR RESUME]
  if resume_population is not None:
      population = resume_population
      if resume_rng_state: rng.setstate(rng_state_from_dict(resume_rng_state))
      loop_start = resume_start_gen + 1
  else:
      population = [Individual(prompt=p, operators="seed") for p in seed_prompts[:5]]

      -- load per-seed checkpoints from previous crashed run
      existing_seeds: dict[int, dict] = {ckpt.seed_idx: ckpt for ckpt in seed_NN_*.json}

      for seed_idx, ind in enumerate(population):
          if seed_idx in existing_seeds:
              restore ind.uid, fitness, failure_examples from checkpoint
              continue
          ind.fitness = compute_fitness(0, ind, problems)   -- 65 Gemini Flash calls
          save seed checkpoint: run_dir/seed_NN_<uid>.json  -- crash recovery

      loop_start = 1

  stagnation_count = 0

  [GENS loop_start .. 12]
  for gen in range(loop_start, GENERATIONS + 1):

      population.sort(by fitness, descending)
      fits = [ind.fitness for ind in population]
      stats = {generation:gen-1, best, mean, worst, std, cv, best_uid, best_prompt}
      if stats.generation > history[-1].generation:
          history.append(stats)

      on_generation(gen-1, population, stats, rng.getstate())
          -> save_checkpoint(run_dir, gen-1, pop, stats, rng_state)

      [CONVERGENCE CHECK]
      cv = pop_std / pop_mean  (if pop_mean > CONVERGENCE_MIN_MEAN=0.01)
      if cv < CONVERGENCE_CV_THRESHOLD=0.01:
          stagnation_count += 1
          if stagnation_count >= CONVERGENCE_PATIENCE=3:
              break   -- EARLY STOP
      else:
          stagnation_count = 0

      [ELITISM -- _elitism frozen elites (0 in default scheme)]
      new_pop = population[:_elitism]

      [BREED remaining slots]
      offspring_idx = 0
      while len(new_pop) < POPULATION_SIZE:
          p1 = _select(population, rng)     -- tournament or roulette
          p2 = _select(population, rng)
          allow_mut = (offspring_idx >= CROSSOVER_ONLY_SLOTS=2)
          child = breed(p1, p2, rng, allow_mutation=allow_mut)
          new_pop.append(child)
          offspring_idx += 1

      [EVALUATE]
      for off_idx, ind in enumerate(new_pop):
          ind.fitness = compute_fitness(gen, ind, problems)

      population = new_pop
      population.sort(by fitness, descending)

  [FINAL STATS]
  population.sort()
  final_stats = {generation:last_completed_gen, best, mean, ...}
  history.append(final_stats)

  return (population[0], history)
\end{lstlisting}

\begin{conceptbox}{Generational replacement with no elites (default)}
Unlike classic elitist GAs, every slot is a new offspring each generation.
The best individual is not automatically preserved.  This increases exploration
at the cost of potentially losing a good solution temporarily.  With POPULATION\_SIZE=5,
even one excellent offspring that emerges from crossover of the best parents is likely
to dominate the next tournament selection.
\end{conceptbox}

\begin{conceptbox}{Three selection schemes (\fn{--scheme} argument)}
The \fn{selection\_scheme} parameter (set via \fn{--scheme} in \fn{main.py}) controls
how parents are chosen each generation:
\begin{itemize}[nosep]
  \item \textbf{tournament} (default): sample 3 at random, take the fittest.
    Strong selection pressure; fast; default for all experiments.
  \item \textbf{roulette}: fitness-proportionate selection.  Individuals with
    higher fitness have a proportionally higher chance of being chosen.
    Falls back to uniform random when all fitnesses are zero.
  \item \textbf{elitism}: same tournament selection plus \fn{elitism\_count=1}
    frozen elite.  The single fittest individual from the previous generation
    is carried over unchanged to guarantee no regression.
\end{itemize}
All three schemes share the same crossover and mutation pipeline.
\end{conceptbox}

\begin{conceptbox}{Per-seed crash recovery}
After each seed is evaluated at gen-0, its result is saved to a file named
\fn{seed\_NN\_<uid>.json}.  If the process crashes during gen-0 (e.g.\ after seed 2 but
before seed 3), the next run finds these checkpoints and skips the already-evaluated
seeds.  This prevents wasting API quota re-evaluating seeds that already succeeded.
\end{conceptbox}


% ==========================================================================
%  CHAPTER 10 --- main.py
% ==========================================================================
\chapter{Orchestration: \pyfile{main.py}}

\section{Purpose}

The entry point.  Handles pre-flight credential checks, dataset loading, run directory
creation, GA execution, holdout evaluation, genealogy report generation, and cost report
saving.  Also implements crash-resume logic: a run can be resumed from any saved
generation checkpoint, with deterministic RNG restoration.

\section{File at a Glance}

\begin{table}[h]
\centering
\begin{tabular}{@{}llp{8cm}@{}}
\toprule
\textbf{Item} & \textbf{Type} & \textbf{Description} \\
\midrule
\fn{save\_checkpoint()}      & function & Serialize full population to generation\_NN.json \\
\fn{\_find\_latest\_run()}   & private  & Find most recently modified run\_* directory \\
\fn{\_load\_resume\_state()} & private  & Load last checkpoint; reconstruct Individual objects \\
\fn{\_holdout\_one()}        & private  & Evaluate one problem against the best prompt \\
\fn{evaluate\_on\_holdout()} & function & Parallel holdout evaluation; returns mean score \\
\fn{main()}                  & function & The 8-step orchestration pipeline \\
\bottomrule
\end{tabular}
\caption{\pyfile{main.py}: all items.}
\end{table}

\section{Function-by-Function Analysis}

\subsection{\fn{save\_checkpoint(run\_dir, gen, population, stats, rng\_state=None)}}

\begin{specbox}{save\_checkpoint(...) -> None}
\textbf{Called by:} the \fn{on\_generation} callback in \fn{main()}\\
\textbf{Output:} \fn{results/run\_*/generation\_NN.json}
\end{specbox}

Saves the full population state (all 5 prompts, fitnesses, parent UIDs, operators,
intermediate steps, failure examples) \emph{and} the MT19937 RNG state:

\begin{lstlisting}[style=example]
data = {
    "generation": gen,
    "stats":      stats,
    "rng_state":  ga.rng_state_to_dict(rng_state),   # or None
    "population": [
        {"uid": ind.uid, "prompt": ind.prompt, "fitness": ind.fitness,
         "parent_a": ind.parent_a, "parent_b": ind.parent_b,
         "operators": ind.operators,
         "intermediate_steps": ind.intermediate_steps,
         "failure_examples": ind.failure_examples}
        for ind in population
    ],
}
gen_file.write_text(json.dumps(data, indent=2))
\end{lstlisting}

\subsection{\fn{\_find\_latest\_run()} and \fn{\_load\_resume\_state()}}

\fn{\_find\_latest\_run()} globs \fn{config.RESULTS\_DIR} for \fn{run\_*} directories
and returns the one with the most recent \fn{st\_mtime}.

\fn{\_load\_resume\_state(run\_dir)} reads all \fn{generation\_NN.json} files, extracts
the last one's population, rebuilds \fn{Individual} objects, and reconstructs the
history list from all checkpoints.  Also reads the saved RNG state for deterministic
continuation.

\begin{lstlisting}[style=chain]
_load_resume_state(run_dir) -> dict:
  checkpoints = sorted(run_dir.glob("generation_*.json"))
  if not checkpoints: raise FileNotFoundError(...)

  history = [data["stats"] for each checkpoint]
  history.sort(key=lambda s: s["generation"])

  last_ckpt = json.loads(checkpoints[-1].read_text())
  population = [
      ga.Individual(prompt=ind["prompt"], fitness=ind["fitness"],
                    uid=ind["uid"], parent_a=ind["parent_a"],
                    parent_b=ind["parent_b"], operators=ind["operators"],
                    intermediate_steps=ind["intermediate_steps"],
                    failure_examples=ind["failure_examples"])
      for ind in last_ckpt["population"]
  ]
  rng_state    = last_ckpt.get("rng_state")
  evolution_ids = json.loads((run_dir/"evolution_ids.json").read_text())

  return {"population": population, "start_gen": last_gen,
          "history": history, "evolution_ids": evolution_ids,
          "rng_state": rng_state}
\end{lstlisting}

\subsection{\fn{evaluate\_on\_holdout(best\_prompt, holdout, run\_dir)} $\to$ \fn{float}}

\begin{specbox}{evaluate\_on\_holdout(best\_prompt, holdout, run\_dir) -> float}
\textbf{Parallelism:} ThreadPoolExecutor(max\_workers=PARALLEL\_EVALS=5)\\
\textbf{Output:} results appended to \fn{holdout\_evaluations.jsonl}
\end{specbox}

Evaluates the best prompt on 200 held-out problems in parallel.  Each call to
\fn{\_holdout\_one()} sleeps a random 0--2 seconds to stagger API requests, then
calls \fn{llm.generate\_code()} and \fn{evaluator.run\_tests\_verbose()}.  Results
are appended to JSONL immediately (thread-safe lock).  Returns the mean score.

\subsection{\fn{main()} --- The 8-Step Pipeline}

\begin{lstlisting}[style=chain]
main():

  [1] Parse args
      --resume [RUN_DIR]   resume a crashed run (AUTO or explicit path)
      --scheme {tournament|roulette|elitism}
                           parent selection scheme (default: tournament)

  [2] Pre-flight: Vertex AI credentials
      check SERVICE_ACCOUNT_PATH exists
      import google.genai, google.oauth2.service_account  (verify installed)
      read service-account.json: log project_id, client_email

  [3] Pre-flight: dataset dirs
      if not APPS_TRAIN.exists() or not APPS_TEST.exists(): sys.exit(1)

  [4] Load dataset
      all_problems = dataset.load_problems()   -- 305 Codeforces problems

  [5] Handle resume vs fresh run
      Fresh:
          evolution_pool, holdout, rest = dataset.split(all_problems, 65, 200, 42)
          run_dir = RESULTS_DIR / "run_YYYYMMDD_HHMMSS_<scheme>"
                  -- e.g. "run_20240501_143022_tournament"
          writes: config.json, evolution_ids.json

      Resume (--resume [path]):
          run_dir = specified path or _find_latest_run()
          if seed checkpoints exist but no generation checkpoints:
              resume_state = None  -- ga.run() handles seed-level resume
          else:
              resume_state = _load_resume_state(run_dir)
          restore evolution_pool from evolution_ids.json

  [6] Set up logging
      tracker.set_log_path(run_dir/"calls.jsonl")
      tracker.set_live_path(run_dir/"costs_live.json")
      tracker.set_costs_live_path(Costs/"cost_<run>.json")
      llm.set_error_log_path(run_dir/"llm_errors.log")
      if resume: tracker.restore_from_jsonl(run_dir/"calls.jsonl")

  [7] Run GA
      best, history = ga.run(
          seed_prompts=SEED_PROMPTS[:POPULATION_SIZE],
          problems=evolution_pool,
          run_dir=run_dir,
          on_generation=lambda gen, pop, stats, rng_s:
              save_checkpoint(run_dir, gen, pop, stats, rng_s),
          resume_population=resume_state["population"] if resume_state else None,
          resume_start_gen= resume_state["start_gen"]  if resume_state else 0,
          resume_history=   resume_state["history"]    if resume_state else None,
          resume_rng_state= resume_state["rng_state"]  if resume_state else None,
          selection_scheme= args.scheme,                  -- "tournament"|"roulette"|"elitism"
          elitism_count=    1 if args.scheme == "elitism" else None,
      )

  [8] Save evolution results
      (run_dir/"history.json").write_text(history)
      (run_dir/"best_prompt.txt").write_text(best.prompt)

  [9] Holdout evaluation
      holdout_fitness = evaluate_on_holdout(best.prompt, holdout, run_dir)
      (run_dir/"summary.json").write_text({elapsed, best_prompt,
          best_fitness_evolution, best_fitness_holdout, best_uid,
          total_generations, resumed})

  [10] Genealogy report
      genealogy.generate_report(run_dir, best.uid)
      -> genealogy_report.txt + genealogy_tree.json

  [11] Cost report
      cost_path = tracker.flush(extra={run_dir, backend, fitnesses, ...})
      -> Code/Costs/cost_TIMESTAMP.json
\end{lstlisting}

\section{Output Files}

\begin{table}[h]
\centering
\begin{tabular}{@{}p{6cm}p{6.5cm}@{}}
\toprule
\textbf{File} & \textbf{Contents} \\
\midrule
\fn{run\_*/config.json}              & Exact hyperparameters snapshot \\
\fn{run\_*/evolution\_ids.json}      & The 65 problem IDs in the evolution pool \\
\fn{run\_*/calls.jsonl}              & Per-call token + cost log (streaming) \\
\fn{run\_*/costs\_live.json}         & Live cost summary (overwritten after every call) \\
\fn{run\_*/llm\_errors.log}          & API retry / error events \\
\fn{run\_*/progress.log}             & Human-readable timestamped timeline \\
\fn{run\_*/evaluations\_genNN.jsonl} & Per-problem full evaluation records \\
\fn{run\_*/generation\_NN.json}      & Full population state + RNG state per generation \\
\fn{run\_*/seed\_NN\_<uid>.json}     & Per-seed checkpoint (gen 0 crash recovery) \\
\fn{run\_*/history.json}             & Best/mean/worst/std/cv per generation \\
\fn{run\_*/best\_prompt.txt}         & The evolved prompt (main artefact) \\
\fn{run\_*/summary.json}             & Final fitnesses (evolution + holdout), runtime \\
\fn{run\_*/holdout\_evaluations.jsonl} & Per-problem holdout evaluation records \\
\fn{run\_*/genealogy\_report.txt}    & Human-readable ancestry chain of the best prompt \\
\fn{run\_*/genealogy\_tree.json}     & Machine-readable full ancestry tree \\
\fn{Costs/cost\_*.json}              & Token counts and estimated cost breakdown \\
\bottomrule
\end{tabular}
\caption{All output files produced by one run.}
\end{table}


% ==========================================================================
%  CHAPTER 11 --- genealogy.py
% ==========================================================================
\chapter{Genealogy: \pyfile{genealogy.py}}

\section{Purpose}

After the GA finishes, \pyfile{genealogy.py} traces the ancestry of the best prompt
backward through all checkpoint files, building a human-readable lineage report and
a machine-readable full ancestry tree.  This allows post-hoc analysis of how the best
prompt evolved: which seed it descended from, which crossover created its key
characteristics, and which mutations refined it.

\section{File at a Glance}

\begin{table}[h]
\centering
\begin{tabular}{@{}llp{8cm}@{}}
\toprule
\textbf{Item} & \textbf{Type} & \textbf{Description} \\
\midrule
\fn{\_load\_all\_individuals()} & private & Build a uid-keyed registry from all generation checkpoints \\
\fn{\_collect\_lineage()}       & private & Walk parent\_a chain from best to seed \\
\fn{\_collect\_full\_tree()}    & private & Recursively collect both parents at every fork \\
\fn{generate\_report()}         & function & Entry point: produce report.txt and tree.json \\
\fn{\_write\_text()}            & private & Simple path.write\_text() wrapper \\
\bottomrule
\end{tabular}
\caption{\pyfile{genealogy.py}: all items.}
\end{table}

\section{Algorithm}

\begin{lstlisting}[style=chain]
generate_report(run_dir, best_uid):

  [1] Load registry
      registry = {}
      for each generation_NN.json checkpoint:
          for each individual in checkpoint["population"]:
              if uid not in registry:
                  registry[uid] = {uid, prompt, fitness, parent_a, parent_b,
                                   operators, generation}

  [2] Primary lineage (parent_a chain only)
      chain = []
      current = best_uid
      while current in registry and current not visited:
          chain.append(registry[current])
          current = registry[current]["parent_a"]

  [3] Write genealogy_report.txt
      For each step in chain:
          print uid, generation, fitness, operators (role: BEST / SEED / ancestor N)
          print full prompt text
          print both parent uids + their fitness and generation

  [4] Full ancestry tree (both parents at every fork)
      def _collect_full_tree(uid, depth=0):
          if depth > 40 or not uid: return {}
          node = {uid, generation, fitness, operators, prompt_preview[:120]}
          if parent_a: node["parent_a"] = _collect_full_tree(parent_a, depth+1)
          if parent_b != parent_a: node["parent_b"] = _collect_full_tree(parent_b, depth+1)
          return node

      tree = _collect_full_tree(best_uid)
      write genealogy_tree.json (JSON, indented)
\end{lstlisting}

\begin{conceptbox}{What the genealogy tells you}
The \fn{genealogy\_report.txt} answers: ``Where did the best prompt come from?''
If the chain is: best $\to$ ancestor 1 $\to$ ancestor 2 $\to$ seed,
you can read exactly which seed it descended from, which crossover blended two parent
strategies, and which inject mutation added the critical strategy.  The \fn{operators}
field at each step says what genetic operation produced that individual.
The \fn{genealogy\_tree.json} shows the full binary tree --- both parents at every fork
--- useful for quantifying how much of the seed population contributed to the final answer.
\end{conceptbox}


% ==========================================================================
%  GLOSSARY
% ==========================================================================
\chapter*{Glossary}
\addcontentsline{toc}{chapter}{Glossary}

\begin{description}[style=nextline, leftmargin=3cm, itemsep=3pt]

\item[API call budget] One full run makes $\approx$4{,}425 Gemini Flash calls (fitness) and
  $\approx$66 Gemini Pro calls (crossover + mutation).  See Chapter~2.

\item[Chromosome] A candidate solution in a GA.  Here: a system prompt string,
  stored as an \fn{Individual} dataclass.

\item[Coefficient of Variation (CV)] $\text{std} / \text{mean}$ of fitnesses in the
  current population.  Used by the convergence detector: CV $<$ 0.01 for 3 consecutive
  generations triggers early stopping.

\item[Convergence detection] Early stopping when the population becomes statistically
  homogeneous (CV below threshold for PATIENCE generations).  Only checked when
  mean fitness $>$ CONVERGENCE\_MIN\_MEAN.

\item[CROSSOVER\_ONLY\_SLOTS] First $N=2$ offspring per generation receive crossover
  only --- no mutation.  Preserves genetic material faithfully while still allowing
  recombination.

\item[Crossover] Genetic operator combining two parents into one offspring.  Here:
  the Optimizer LLM (Gemini 2.5 Pro) semantically blends two parent prompts.
  Fires with probability 0.8.

\item[Early stopping (test)] In \fn{run\_tests\_verbose()}: halting the test loop after
  EARLY\_STOP\_FAILURES=3 consecutive failures.  Score denominator is still total tests.

\item[Elitism] Copying the top $k$ individuals unchanged.  \fn{ELITISM\_COUNT=0} in
  the current design --- no frozen elites.

\item[Evolution pool] The fixed 65 Codeforces problems used as the fitness benchmark,
  drawn by stratified sampling (seed 42).

\item[Failure examples] Up to 5 problem-level failure records stored in each
  \fn{Individual.failure\_examples}.  Passed to \fn{mutate\_inject()} to focus new
  strategy additions on actual failure modes.

\item[Fitness] Mean partial-credit score across 65 problems.  Range: [$-1.0$, 1.0]
  (infra failures score $-1.0$; otherwise [0.0, 1.0]).

\item[Gemini Flash] \fn{gemini-2.5-flash} --- the Target model.  Used for code
  generation (fitness evaluation).  Temperature 0.0, thinking\_budget 512.

\item[Gemini Pro] \fn{gemini-2.5-pro} --- the Optimizer model.  Used for crossover
  and mutation.  Temperature 0.7, thinking\_budget 2048.

\item[Genealogy] Ancestry report tracing the best evolved prompt back to its seed.
  Produced by \pyfile{genealogy.py} after the GA completes.

\item[Generational replacement] Every slot is filled by a new offspring each generation
  (no frozen elites).  Creates strong selection pressure via tournament selection.

\item[Held-out set] 200 problems never seen during evolution.  Used to test
  generalisation of the final evolved prompt.

\item[Individual] One GA population member: a prompt with fitness, uid, parent uids,
  operator history, intermediate steps, and failure examples.

\item[Intermediate steps] \fn{Individual.intermediate\_steps}: dict recording the
  prompt text after each genetic operator (\fn{after\_crossover}, \fn{after\_inject},
  etc.).  Stored in checkpoints for post-hoc analysis.

\item[Meta-prompt] A prompt sent to the Optimizer LLM instructing it to modify another
  prompt (used in crossover and all three mutation operators).

\item[Mutation] Modifying a single chromosome.  Three independent types: inject (25\%),
  delete (15\%), rephrase (10\%).  Only fires on offspring with index $\geq$ CROSSOVER\_ONLY\_SLOTS.

\item[Partial credit] Score = (tests passed) / (total tests).  Provides a smooth
  gradient vs.\ binary pass/fail.

\item[PARALLEL\_EVALS] Number of problems evaluated simultaneously per individual
  (5, using a ThreadPoolExecutor).

\item[Resume] Running with \fn{--resume} loads the last generation checkpoint,
  restores the population and RNG state, and continues from the next generation.
  Seed-level checkpoints (\fn{seed\_NN\_*.json}) also allow gen-0 crash recovery.

\item[RNG state] The full MT19937 internal state serialized via
  \fn{ga.rng\_state\_to\_dict()}.  Saved in each generation checkpoint; restored on
  resume to make tournament selection decisions deterministic.

\item[Service account] A Google Cloud IAM service account JSON credential file
  (\fn{service-account.json}).  Used by \fn{llm.\_get\_api\_client()} to authenticate
  with Vertex AI without interactive login.

\item[Singleton] The \fn{tracker = CostTracker()} instance at module level; the
  \fn{\_api\_client} Vertex AI client.  Both shared across all threads.

\item[Stratified split] The evolution pool is built by drawing a fixed number of
  problems from each difficulty tier (800--1700) plus all harder problems (1800/1900),
  ensuring a consistent difficulty distribution.

\item[Thinking tokens] Gemini 2.5 Flash/Pro internal chain-of-thought tokens.
  Billed at the output token rate.  Tracked separately via
  \fn{usage\_metadata.thoughts\_token\_count} and added to output token count.

\item[Tournament selection] Sample $k=3$ individuals at random; return the fittest.
  Balances exploitation with exploration.

\item[Two-model architecture] Flash (code generation, target) is optimised by the GA;
  Pro (reasoning, creative) performs the genetic operators.

\item[Vertex AI] Google Cloud's managed AI platform.  LLM calls use the
  \fn{google-genai} Python SDK connected to the \fn{us-central1} endpoint with
  service account credentials.

\end{description}

\end{document}
